\documentclass[11pt,reqno]{amsart}
\usepackage[utf8]{inputenc}
\usepackage[T1]{fontenc}

\usepackage[margin=1.15in]{geometry}
\usepackage{amsmath,amssymb,amsthm,mathtools}
\usepackage{tikz-cd}
\usepackage{enumitem}
\usepackage{booktabs}
\usepackage[colorlinks=true,linkcolor=blue!50!black,citecolor=green!40!black,urlcolor=blue!50!black]{hyperref}
\usepackage[capitalise,nameinlink,noabbrev]{cleveref}

\setlist[enumerate]{label=\textup{(\arabic*)}}
\setlength{\emergencystretch}{3em}

% ---------------------------------------------------------------
% Theorem environments
% ---------------------------------------------------------------
\theoremstyle{plain}
\newtheorem{theorem}{Theorem}[section]
\newtheorem{proposition}[theorem]{Proposition}
\newtheorem{lemma}[theorem]{Lemma}
\newtheorem{corollary}[theorem]{Corollary}
\newtheorem{conjecture}[theorem]{Conjecture}
\newtheorem{question}[theorem]{Question}
\newtheorem{maintheoreminner}{Theorem}
\newenvironment{maintheorem}[1]{\renewcommand{\themaintheoreminner}{#1}\begin{maintheoreminner}}{\end{maintheoreminner}}
\theoremstyle{definition}
\newtheorem{definition}[theorem]{Definition}
\newtheorem{example}[theorem]{Example}
\newtheorem{setting}[theorem]{Setting}
\newtheorem{notation}[theorem]{Notation}
\theoremstyle{remark}
\newtheorem{remark}[theorem]{Remark}
\newtheorem*{intuition}{Intuition}

\crefname{maintheoreminner}{Theorem}{Theorems}
\crefname{setting}{Setting}{Settings}
\crefname{notation}{Notation}{Notations}
\crefname{question}{Question}{Questions}
\crefname{conjecture}{Conjecture}{Conjectures}

% ---------------------------------------------------------------
% Macros
% ---------------------------------------------------------------
\newcommand{\kk}{\Bbbk}
\DeclareMathOperator{\Hom}{Hom}
\DeclareMathOperator{\End}{End}
\DeclareMathOperator{\Ext}{Ext}
\DeclareMathOperator{\Filt}{Filt}
\DeclareMathOperator{\add}{add}
\DeclareMathOperator{\real}{real}
\DeclareMathOperator{\gldim}{gl.dim}
\DeclareMathOperator{\pd}{pd}
\DeclareMathOperator{\soc}{soc}
\DeclareMathOperator{\rad}{rad}
\DeclareMathOperator{\tr}{tr}
\DeclareMathOperator{\im}{Im}
\DeclareMathOperator{\cone}{cone}
\DeclareMathOperator{\RHom}{\mathbf{R}Hom}
\newcommand{\modd}{\operatorname{mod}}
\newcommand{\proj}{\operatorname{proj}}
\newcommand{\Db}{\mathcal{D}^b}
\newcommand{\Kb}{K^b}
\newcommand{\thick}[2][\cD]{\operatorname{thick}_{#1}(#2)}
\newcommand{\Pn}[2]{\mathbb{P}^{#1}[#2]}
\newcommand{\nablabar}{\overline{\nabla}}
\newcommand{\Deltabar}{\overline{\Delta}}
\newcommand{\cD}{\mathcal{D}}
\newcommand{\cA}{\mathcal{A}}
\newcommand{\cH}{\mathcal{H}}
\newcommand{\cT}{\mathcal{T}}
\newcommand{\cC}{\mathcal{C}}
\newcommand{\cF}{\mathcal{F}}
\newcommand{\cN}{\mathcal{N}}
\newcommand{\cM}{\mathcal{M}}
\newcommand{\cQ}{\mathcal{Q}}
\newcommand{\cU}{\mathcal{U}}
\newcommand{\cV}{\mathcal{V}}
\newcommand{\cO}{\mathcal{O}}
\newcommand{\pH}{{}^{p}\!H}

\title[Local-exceptional sequences and stratified algebras]{A stratified Bodzenta--Bondal criterion: local-exceptional sequences, tilting objects and stratified algebras}
\date{October 2026}

\begin{document}

\begin{abstract}
Bodzenta and Bondal characterised highest weight categories in terms of a full exceptional sequence and its left dual. We prove a stratified analogue, formulated entirely in terms of objects of a triangulated category $\cD$:
\begin{itemize}
\item the exceptional sequence is replaced by a sequence of \emph{local-exceptional} objects $E_s$, which have no self-extensions in non-zero degrees but may have arbitrary local endomorphism algebras;
\item the left dual is replaced by a \emph{Hom-dual} sequence of ``simple-like'' objects $N_s$, which generates $\cD$, together with its right dual in the sense of mutations.
\end{itemize}
We show the following. Suppose neither sequence has negative extensions. Then the iterated universal extension $G$ of the $E_s$ is a tilting object of the subcategory they generate, its endomorphism algebra $\Lambda$ is standardly stratified, and $\cD\simeq\Db(\modd\Lambda)$; under this equivalence the $E_s$ and $N_s$ become the standard and proper costandard modules. The proof is an induction on the number of strata, modelled on heredity chains. t-structures appear only inside the proof, as the t-structure determined by the given objects.

A version with Hom-duality only in low degrees describes hearts that are standardly stratified although $\cD$ is not their derived category. Conversely, the standard modules of a standardly stratified algebra form a local-exceptional sequence whose Hom-dual partners are not perfect unless the algebra is quasi-hereditary. As an application we discuss sequences of $\Pn nk$-like objects.
\end{abstract}

\maketitle

\setcounter{tocdepth}{1}
\tableofcontents

% ===============================================================
\section{Introduction}
% ===============================================================

\subsection{Background}
Highest weight categories, introduced by Cline, Parshall and Scott \cite{CPS88}, are precisely the module categories of quasi-hereditary algebras. Their bounded derived categories carry two distinguished full exceptional sequences:
\begin{itemize}
\item the standard objects $(\Delta(\lambda_1),\dots,\Delta(\lambda_n))$;
\item the costandard objects, which form the left dual sequence $(\nabla(\lambda_n),\dots,\nabla(\lambda_1))$.
\end{itemize}
Bodzenta and Bondal \cite{BB26} proved that this characterises highest weight categories. The heart of the t-structure glued along a full exceptional sequence is highest weight, with the given sequence and its left dual as standard and costandard objects, if and only if neither sequence has negative extensions (\cref{thm:BB-3.1}). Moreover, highest weight structures on an abelian category $\cA$ correspond bijectively to pairs of dual full exceptional sequences in $\Db(\cA)$ consisting of objects of $\cA$ (\cref{thm:BB-3.8}).

\begin{intuition}
A full exceptional sequence is the same as an admissible filtration of $\cD$ whose subquotients (the \emph{strata}) are all equivalent to $\Db(\kk)$. In $\Db(\kk)$ the object $\kk$ is simultaneously the simple object and the projective object of the heart. The standard object $E_s$ is its extension by zero $j_{s!}(\kk)$, and the costandard object $F_s$ is its extension $j_{s*}(\kk)$.
\end{intuition}

\subsection{The question}
\emph{Standardly stratified} and \emph{properly stratified} algebras \cite{CPS96,Dla96,AHLU00} generalise quasi-hereditary algebras. The endomorphism algebras $B_i=\End(\Delta(i))$ of the standard modules are allowed to be arbitrary local algebras instead of $\kk$. We address the following question.

\begin{question}\label{q:main}
Is there a generalisation of dual pairs of exceptional sequences which recovers standardly stratified algebras in the same way as dual pairs of exceptional sequences recover quasi-hereditary algebras?
\end{question}

The question arose from a more specific one, concerning $\Pn nk$-like objects \cite{HT06,Kru18,HK19}. These are objects $P$ with $\End^\bullet(P)\cong\kk[t]/(t^{n+1})$, $\deg t=k$, which generalise exceptional objects (the case $n=0$) in a different direction. They are treated as an application in \cref{sec:P-objects}.

The natural candidates for standard objects are \emph{local-exceptional} objects (\cref{def:local-exc}). These are objects $E$ with $\Hom(E,E[l])=0$ for $l\ne0$ and $\End(E)$ local. Two observations show that passing from exceptional to local-exceptional sequences is not formal.
\begin{itemize}
\item The category generated by a local-exceptional object $E$ is $\thick{E}\simeq\Kb(\proj B)$, where $B=\End(E)$. It contains no object playing the role of the simple $B$-module (\cref{prop:local-exc}). In important cases it has no non-trivial t-structure at all \cite{HJY13,Bonf26}.
\item The left dual of a local-exceptional sequence pairs with it via $\Hom^\bullet(E_s,F_s)\cong B_s$ rather than $\kk$. The correct partners satisfy $\Hom^\bullet(E_s,N_s)\cong\kk$. For stratified algebras these are the proper costandard modules, which lie outside the subcategory generated by the $E_s$ unless $B_s=\kk$ (\cref{mthm:C}).
\end{itemize}
Consequently the dual sequence must be part of the data, and the ambient category must be generated by it.

\subsection{Main results}
All hypotheses below concern objects of a triangulated category $\cD$ and morphisms between them; no t-structure on $\cD$ is assumed. Subcategories are always formed in a named ambient category; for instance $\thick{S}$ denotes the thick subcategory of $\cD$ generated by $S$ (\cref{sec:conventions}).

Let $(E_1,\dots,E_m)$ be a local-exceptional sequence, i.e.\ each $E_s$ is local-exceptional and $\Hom^\bullet(E_t,E_s)=0$ for $t>s$. A \emph{Hom-dual triple} $(E_\bullet,\Deltabar_\bullet,N_\bullet)$ (\cref{def:hom-dual}) consists of this sequence, objects $N_1,\dots,N_m$ and $\Deltabar_1,\dots,\Deltabar_m$, and morphisms $\delta_s\colon\Deltabar_s\to N_s$, satisfying the following four conditions.
\begin{itemize}
\item[\textup{(HD1)}] $\Hom^\bullet(E_s,N_t)\cong\delta_{st}\kk$, concentrated in degree $0$.
\item[\textup{(HD2)}] $\Hom^\bullet(N_s,N_t)=0$ for $s<t$. Moreover $\End(N_s)=\kk$, and $\Hom(N_s,N_s[l])=0$ for $l<0$.
\item[\textup{(HD3)}] $\cD=\thick{N_1,\dots,N_m}$.
\item[\textup{(HD4)}] $\cone(\delta_s)\in\thick{N_1,\dots,N_{s-1}}$, and $\Hom^\bullet(\Deltabar_s,N_t)=0$ for $t<s$.
\end{itemize}
Condition (HD4) says that $(\Deltabar_1,\dots,\Deltabar_m)$ is the right dual of $(N_m,\dots,N_1)$ in the sense of mutations. Equivalently, each $\thick{N_1,\dots,N_{s-1}}$ is left admissible in $\thick{N_1,\dots,N_s}$ (\cref{rem:HD4}). The objects $\Deltabar_s$ play the role of proper standard objects.

\begin{maintheorem}{B}[\cref{thm:B}]\label{mthm:B}
Let $\cD$ be an algebraic triangulated category, and let $(E_\bullet,\Deltabar_\bullet,N_\bullet)$ be a Hom-dual triple in $\cD$. Let $G$ be the iterated universal extension of $(E_1,\dots,E_m)$ in $\cD$, and put $\Lambda=\End_\cD(G)$. The following are equivalent:
\begin{enumerate}
\item $\Hom(E_s,E_t[l])=0$ and $\Hom(N_s,N_t[l])=0$ for all $s,t$ and all $l<0$;
\item $G$ is a tilting object of $\thick{E_1,\dots,E_m}$, the algebra $\Lambda$ is standardly stratified, and there is a triangle equivalence $\cD\simeq\Db(\modd\Lambda)$ sending $G$ to $\Lambda$, $E_s$ to the standard module $\Delta(s)$, and $N_s$ to the proper costandard module $\nablabar(s)$.
\end{enumerate}
In this case:
\begin{itemize}
\item $\End(\Delta(s))\cong\End(E_s)$;
\item $\Lambda$ is quasi-hereditary if and only if all $E_s$ are exceptional;
\item $\thick{E_1,\dots,E_m}\simeq\Kb(\proj\Lambda)$, and this is a proper subcategory of $\cD$ unless $\Lambda$ is quasi-hereditary.
\end{itemize}
\end{maintheorem}

For an exceptional sequence, the left dual is a Hom-dual sequence whose right dual is the sequence itself (\cref{prop:exc-dual}). In that case $\cD=\thick{E_1,\dots,E_n}$, and \cref{mthm:B} becomes \cref{thm:BB-3.1}.

A t-structure appears only in the proof. It is the t-structure on $\cD$ cut out by the objects $E_s$ and $N_s$ (\cref{prop:triple-t}), which a posteriori is the $G$-tilting t-structure (\cref{rem:G-tilting}).

\Cref{mthm:B} is a special case of a statement about \emph{stratifying triples} (\cref{def:triple}). In a stratifying triple:
\begin{itemize}
\item the Hom-duality (HD1) is only required in degrees $\le1$;
\item the objects $\Delta_s$ playing the role of the $E_s$ need not be local-exceptional;
\item the conclusion concerns the heart of the t-structure cut out by the objects, rather than $\cD$ itself.
\end{itemize}

\begin{maintheorem}{A}[\cref{thm:A}]\label{mthm:A}
Let $(\Delta_\bullet,\Deltabar_\bullet,N_\bullet)$ be a stratifying triple in $\cD$ with Hom-duality up to degree $2$, i.e.\ $\Hom(\Delta_s,N_s[l])=0$ for $l\in\{1,2\}$. Suppose that $\Hom(\Delta_s,\Delta_t[l])=0=\Hom(N_s,N_t[l])$ for all $s,t$ and $l<0$. Consider the t-structure
\[\cD^{\le0}=\{X:\Hom(X,N_s[l])=0\ \forall s,\ l<0\},\qquad\cD^{\ge0}=\{X:\Hom(\Delta_s,X[l])=0\ \forall s,\ l<0\},\]
with heart $\cA$. Then:
\begin{enumerate}
\item The iterated universal extension $G$ of $(\Delta_1,\dots,\Delta_m)$ is a projective generator of $\cA$.
\item $\cA\simeq\modd\End(G)$ is standardly stratified, with standard objects $\Delta_s$ and proper costandard objects $N_s$.
\item If $\cD$ is algebraic, then $\cD\simeq\Db(\cA)$ via realisation if and only if Hom-duality holds in all degrees.
\end{enumerate}
\end{maintheorem}

\begin{maintheorem}{C}[\cref{thm:C}]\label{mthm:C}
Let $\Lambda$ be a standardly stratified algebra, and put $B_i=\End(\Delta(i))$.
\begin{enumerate}
\item The standard modules form a full local-exceptional sequence in $\Kb(\proj\Lambda)$ with endomorphism algebras $B_i$.
\item $\Hom^\bullet(\Delta(i),\nablabar(i))\cong\kk$, and $\Hom^\bullet(\Delta(j),\nablabar(i))=0$ for $j>i$.
\item If $B_i\neq\kk$, then $\nablabar(i)\notin\Kb(\proj\Lambda)$.
\item $\gldim\Lambda<\infty$ if and only if $\Lambda$ is quasi-hereditary.
\end{enumerate}
\end{maintheorem}

\Cref{cor:BB38} is the analogue of \cref{thm:BB-3.8}. It is a bijection between:
\begin{itemize}
\item \emph{dualisable} standardly stratified structures, i.e.\ those whose standard and proper costandard objects extend to a Hom-dual triple; and
\item pairs of sequences of objects of $\cA$ that extend to Hom-dual triples.
\end{itemize}
Quasi-hereditary structures are dualisable. A standardly stratified structure is dualisable, with the proper standard modules as right dual, as soon as two conditions hold: injective modules have $\nablabar$-filtrations, and $\Ext^\bullet(\Deltabar(s),\nablabar(t))=0$ for $t<s$ (\cref{prop:dualisable}). Whether every standardly stratified structure is dualisable is open (\cref{q:HD4}).

\subsection{Method: induction on strata}
As in the theory of quasi-hereditary algebras, all proofs are by induction, peeling off one stratum at a time.
\begin{itemize}
\item \emph{Abelian side.} \Cref{lem:strat-induction} characterises standardly stratified categories inductively. Let $P(m)$ be the projective cover of the maximal simple object. The trace of $P(m)$ in every projective object must be a direct sum of copies of $P(m)$ (a ``standardly stratifying'' ideal), and the quotient category must again be standardly stratified. For quasi-hereditary algebras this is the characterisation by heredity chains (\cref{rem:heredity}).
\item \emph{Peeling off a stratum.} By \cref{lem:peel}, restricting a stratifying triple to the thick subcategory of $\cD$ generated by $N_1,\dots,N_{m-1}$ gives a stratifying triple. The last stratum is described by a \emph{local pair} $(\Pi,S)$ (\cref{def:local-pair}): a simple-minded object $S$ together with an object $\Pi$, which turns out to be its projective cover.
\item \emph{Induction step} (\cref{prop:step}). The universal extensions by $\Delta_m$ of the previous projective generator form a projective generator. The trace condition of \cref{lem:strat-induction} then follows formally from the universal-extension triangles. The only non-formal input is Ext-projectivity in degree $2$, which is where Hom-duality in degree $2$ enters.
\item \emph{One gluing lemma.} \Cref{prop:triple-t} constructs the t-structure determined by the objects; this is the only place where BBD gluing is used. One application of the realisation functor then gives $\cD\simeq\Db(\modd\Lambda)$.
\end{itemize}

\subsection{Application: \texorpdfstring{$\mathbb{P}^n[k]$}{Pn[k]}-like objects}
In \cref{sec:P-objects} we specialise to $\Pn nk$-like objects, using Bonfert's formality and classification results \cite{Bonf26}.
\begin{itemize}
\item \emph{$k=0$.} $\Pn n0$-like objects are exactly the local-exceptional objects with endomorphism algebra $\kk[t]/(t^{n+1})$, so \cref{mthm:B,mthm:C} apply.
\item \emph{$k\ge1$.} Let $(P_1,\dots,P_m)$ be a semi-orthogonal sequence of $\Pn{n_s}{k_s}$-like objects with a left dual $(F_m,\dots,F_1)$. It gives a stratifying triple $(\Delta_\bullet,P_\bullet,F_\bullet)$, in which the $P_s$ play the role of proper standard objects (\cref{lem:P-triple}). \Cref{mthm:A} then yields:
\begin{itemize}
\item highest weight hearts when all $k_s\ge3$;
\item standardly stratified hearts with dual-number endomorphism algebras when $k_s=1$ and $n_s\ge3$;
\item hearts that are not module categories when $k_s=n_s=1$.
\end{itemize}
In all cases with $n_s\ge1$, the ambient category is not the derived category of the heart (\cref{thm:Pseq}). The case $k=2$, which contains the Huybrechts--Thomas $\mathbb{P}^n$-objects, is not decided.
\end{itemize}

\subsection*{Sources and status of the statements}
Every statement attributed to the literature is quoted with its statement, in our notation. All statements without a citation are proved here, except those labelled \emph{Question}. From Bonfert's preprint \cite{Bonf26} we use the formality and classification results quoted in \cref{sec:P-objects}. Our computation of the hearts of the categories generated by $\Pn nk$-like objects (\cref{prop:heart-P}) overlaps with \cite[Lemmas~3.11 and~4.15]{Bonf26}, and the argument of \cref{lem:recognise-proj} is that of \cite[Lemma~3.12]{Bonf26}. We include proofs in the form needed here.

\subsection*{Organisation}
\begin{itemize}
\item \Cref{sec:prelim} collects definitions and the results we use.
\item \Cref{sec:local} introduces local-exceptional objects and local pairs.
\item \Cref{sec:triples} introduces stratifying triples, and constructs the associated filtration and t-structure.
\item \Cref{sec:gluing} contains the induction step and \cref{mthm:A}.
\item \Cref{sec:criterion} proves the stratified Bodzenta--Bondal criterion (\cref{mthm:B}).
\item \Cref{sec:stratified} treats standardly stratified algebras (\cref{mthm:C}) and the converse.
\item \Cref{sec:P-objects} contains the application to $\Pn nk$-like objects.
\item \Cref{sec:questions} lists open questions.
\end{itemize}

% ===============================================================
\section{Preliminaries}\label{sec:prelim}
% ===============================================================

\subsection{Conventions}\label{sec:conventions}
Throughout, $\kk$ is an algebraically closed field. Triangulated categories are $\kk$-linear, \emph{Hom-finite} (all $\Hom(X,Y)$ are finite-dimensional) and idempotent complete; in particular they are Krull--Schmidt. A triangulated category is \emph{algebraic} if it is equivalent to the homotopy category $H^0$ of a pretriangulated dg category (equivalently, to the stable category of a Frobenius exact category). For objects $X,Y$ we write $\Hom^l(X,Y):=\Hom(X,Y[l])$ and $\Hom^\bullet(X,Y):=\bigoplus_{l\in\mathbb{Z}}\Hom^l(X,Y)$, so that $\End^\bullet(X)$ is a graded algebra.

\emph{Subcategories and ambient categories.} Every subcategory is formed inside an ambient triangulated category, which is always named. For a class $\mathcal S$ of objects of $\cD$:
\begin{itemize}
\item $\thick[\cD]{\mathcal S}$ is the smallest thick subcategory (triangulated and closed under direct summands) of $\cD$ containing $\mathcal S$. We write $\thick{\mathcal S}$ when $\cD$ is the ambient category of the statement at hand.
\item $\Filt_\cD(\mathcal S)$ is the smallest full subcategory of $\cD$ containing $0$ and $\mathcal S$ and closed under extensions in $\cD$. We often omit the subscript.
\item $\mathcal S^\perp=\{Y\in\cD:\Hom^\bullet(S,Y)=0\ \forall S\in\mathcal S\}$, and ${}^\perp\mathcal S$ is defined dually; both are taken in the ambient category $\cD$.
\end{itemize}
If $\cD'\subset\cD$ is a thick subcategory containing $\mathcal S$, then $\thick[\cD']{\mathcal S}=\thick[\cD]{\mathcal S}$ and $\Filt_{\cD'}(\mathcal S)=\Filt_\cD(\mathcal S)$. For instance, for $\mathcal S\subset\Kb(\proj\Lambda)$ the thick closure is the same in $\Kb(\proj\Lambda)$ and in $\Db(\modd\Lambda)$. We write $\simeq$ for equivalences of abstract triangulated categories, as in $\thick[\cD]{E}\simeq\Kb(\proj B)$; such a statement does not identify the two sides as subcategories of a common category.

Modules are finite-dimensional right modules; $\modd\Lambda$ denotes their category, and $\Kb(\proj\Lambda)\subset\Db(\modd\Lambda)$ the subcategory of perfect complexes. In an abelian length category, $[X:L]$ denotes the composition multiplicity of a simple object $L$ in $X$. For a set $\Theta$ of objects, $\cF(\Theta)$ is the full subcategory of objects admitting a finite filtration with subquotients isomorphic to objects of $\Theta$.

\subsection{t-structures and hearts}

\begin{definition}[{\cite[D\'ef.~1.3.1]{BBD82}}]
A \emph{t-structure} on $\cD$ is a pair $(\cD^{\le0},\cD^{\ge0})$ of strictly full subcategories such that, with $\cD^{\le n}=\cD^{\le0}[-n]$ and $\cD^{\ge n}=\cD^{\ge0}[-n]$:
\begin{enumerate}
\item $\cD^{\le0}[1]\subset\cD^{\le0}$ and $\cD^{\ge0}[-1]\subset\cD^{\ge0}$;
\item $\Hom(X,Y)=0$ for $X\in\cD^{\le0}$ and $Y\in\cD^{\ge1}$;
\item every $X$ fits into a triangle $X'\to X\to X''\to X'[1]$ with $X'\in\cD^{\le0}$ and $X''\in\cD^{\ge1}$.
\end{enumerate}
The \emph{heart} $\cA=\cD^{\le0}\cap\cD^{\ge0}$ is an abelian category, with cohomology functors $H^l\colon\cD\to\cA$. The t-structure is \emph{bounded} if every object lies in $\cD^{\le a}\cap\cD^{\ge b}$ for some $a,b$. It has \emph{length heart} if $\cA$ is a finite length category with finitely many isomorphism classes of simple objects.
\end{definition}

Short exact sequences in $\cA$ are exactly the triangles in $\cD$ all of whose vertices lie in $\cA$. Consequently $\cA$ is closed under extensions in $\cD$.

\begin{proposition}[{\cite[Rem.~3.1.17]{BBD82}}]\label{prop:ext-comparison}
Let $\cA$ be the heart of a t-structure on $\cD$, and let $X,Y\in\cA$. The natural maps $\Ext^l_\cA(X,Y)\to\Hom_\cD(X,Y[l])$, which send a Yoneda extension to the composite of the corresponding connecting morphisms, are bijective for $l\le1$ and injective for $l=2$. They are compatible with Yoneda product and composition.
\end{proposition}

\begin{theorem}[{Realisation; \cite[\S3.1]{BBD82}, \cite{Bei87}, \cite{Let24}}]\label{thm:realisation}
Let $\cD$ be algebraic (more generally, the base of a filtered triangulated category), and let $\cA$ be the heart of a bounded t-structure on $\cD$.
\begin{enumerate}
\item There is an exact functor $\real\colon\Db(\cA)\to\cD$ which restricts to the inclusion on $\cA$ and induces the maps of \cref{prop:ext-comparison} on extension groups.
\item $\real$ is fully faithful if and only if the following \emph{effaceability} condition holds: for all $X,Y\in\cA$, all $n\ge2$ and all $f\in\Hom_\cD(X,Y[n])$, there is an epimorphism $X'\twoheadrightarrow X$ in $\cA$ such that the composite $X'\to X\xrightarrow{f}Y[n]$ vanishes.
\item In this case $\real$ is an equivalence, since $\cD$ is generated by $\cA$ under extensions and shifts.
\end{enumerate}
\end{theorem}

\begin{remark}
In \cite[Prop.~3.1.16]{BBD82} the effaceability condition is formulated with monomorphisms $Y\hookrightarrow Y'$. The version with epimorphisms follows by applying the statement to $\cD^{\mathrm{op}}$. The existence of $\real$ for algebraic triangulated categories is established, for instance, in \cite{Let24}.
\end{remark}

\subsection{Simple-minded collections, silting objects and co-t-structures}

\begin{definition}\label{def:smc}
A finite set $\{S_1,\dots,S_r\}$ of objects of $\cD$ is a \emph{simple-minded collection} if
\begin{enumerate}
\item[(SM1)] $\Hom(S_i,S_j[l])=0$ for all $i,j$ and all $l<0$;
\item[(SM2)] $\Hom(S_i,S_j)=\delta_{ij}\kk$;
\item[(SM3)] $\thick{S_1,\dots,S_r}=\cD$.
\end{enumerate}
An object $T$ is \emph{silting} if $\Hom(T,T[l])=0$ for all $l>0$ and $\thick{T}=\cD$. It is \emph{tilting} if moreover $\Hom(T,T[l])=0$ for all $l<0$.
\end{definition}

\begin{intuition}
Simple-minded collections behave like the simple modules of a finite-dimensional algebra, and silting objects behave like its projective modules. The former determine t-structures, the latter co-t-structures (weight structures). For $\Db(\modd\Lambda)$ the two notions correspond to each other under the bijections of Koenig and Yang (\cref{thm:KY}). This ``simple versus projective'' dichotomy is the mechanism behind \cref{thm:HJY} and behind all results below.
\end{intuition}

\begin{theorem}[{\cite{AN09}; see also \cite[\S5]{KY14}}]\label{thm:smc}
Let $\{S_1,\dots,S_r\}$ be a simple-minded collection in a Hom-finite Krull--Schmidt triangulated category $\cD$. Put
\[\cD^{\le0}:=\Filt\{S_i[l]: 1\le i\le r,\ l\ge0\},\qquad \cD^{\ge0}:=\Filt\{S_i[l]:1\le i\le r,\ l\le0\}.\]
Then $(\cD^{\le0},\cD^{\ge0})$ is a bounded t-structure on $\cD$. Its heart is $\Filt\{S_1,\dots,S_r\}$, a length category whose simple objects are exactly $S_1,\dots,S_r$.
\end{theorem}

\begin{definition}[{\cite{Bon10,Pau08}}]
A \emph{co-t-structure} (or \emph{weight structure}) on $\cD$ is a pair $(\cU,\cV)$ of full subcategories with the following properties: both are closed under direct summands; $\cU[-1]\subset\cU$ and $\cV[1]\subset\cV$; $\Hom(\cU,\cV)=0$; and every object $X$ fits into a triangle $U\to X\to V\to U[1]$ with $U\in\cU$ and $V\in\cV$. Its \emph{co-heart} is $\cU\cap\cV[-1]$. It is \emph{bounded} if $\bigcup_l\cU[l]=\cD=\bigcup_l\cV[l]$.
\end{definition}

\begin{theorem}[{\cite{Bon10,AI12}}]\label{thm:silting-cot}
Let $T$ be a silting object of $\cD$. Let $\cU$ (resp.\ $\cV$) be the smallest full subcategory closed under extensions and direct summands containing all $T[l]$ with $l\le0$ (resp.\ $l\ge1$). Then $(\cU,\cV)$ is a bounded co-t-structure with co-heart $\add T$.
\end{theorem}

\begin{theorem}[{Koenig--Yang \cite[Thm.~6.1]{KY14}}]\label{thm:KY}
For a finite-dimensional algebra $\Lambda$ there are bijections between:
\begin{enumerate}
\item equivalence classes of silting objects in $\Kb(\proj\Lambda)$;
\item equivalence classes of simple-minded collections in $\Db(\modd\Lambda)$;
\item bounded t-structures on $\Db(\modd\Lambda)$ with length heart;
\item bounded co-t-structures on $\Kb(\proj\Lambda)$.
\end{enumerate}
\end{theorem}

\begin{theorem}[{Keller \cite{Kel94}}]\label{thm:keller}
Let $\cD$ be an algebraic idempotent complete triangulated category, and let $T\in\cD$ be a tilting object. Then there is a triangle equivalence $\cD\simeq\Kb(\proj\End(T))$ sending $T$ to the free module of rank one.
\end{theorem}

\subsection{Recollements and gluing}

\begin{definition}\label{def:recollement}
A full triangulated subcategory $\cN\subset\cM$ is \emph{admissible} if the inclusion $i_*\colon\cN\to\cM$ has a left adjoint $i^*$ and a right adjoint $i^!$. Put $\cQ=\cM/\cN$ (Verdier quotient), and let $j^*\colon\cM\to\cQ$ be the quotient functor. It has a left adjoint $j_!$ and a right adjoint $j_*$, both fully faithful, which identify $\cQ$ with ${}^\perp\cN$ and with $\cN^\perp$ respectively. We obtain a \emph{recollement}
\[
\begin{tikzcd}[column sep=large]
\cN \arrow[r,"i_*" description] & \cM \arrow[l,bend right=40,"i^*"'] \arrow[l,bend left=40,"i^!"] \arrow[r,"j^*" description] & \cQ \arrow[l,bend right=40,"j_!"'] \arrow[l,bend left=40,"j_*"]
\end{tikzcd}
\]
with $i^*j_!=0$, $i^!j_*=0$, $j^*i_*=0$, $j^*j_!\cong\mathrm{id}\cong j^*j_*$, and functorial triangles
\[ j_!j^*X\to X\to i_*i^*X\to,\qquad i_*i^!X\to X\to j_*j^*X\to.\]
Applying $i^!$ to the first triangle for $X=j_*Y$ gives $i^!j_!Y\cong i^*j_*Y[-1]$.
\end{definition}

\begin{theorem}[{Gluing \cite[Thm.~1.4.10 and \S1.4]{BBD82}}]\label{thm:gluing}
In a recollement as above, let $\cN$ and $\cQ$ carry t-structures. Then
\begin{align*}
\cM^{\le0}&=\{X: j^*X\in\cQ^{\le0},\ i^*X\in\cN^{\le0}\},\\
\cM^{\ge0}&=\{X: j^*X\in\cQ^{\ge0},\ i^!X\in\cN^{\ge0}\}
\end{align*}
is a t-structure on $\cM$, the \emph{glued t-structure}. With respect to it, $i_*$ and $j^*$ are t-exact, $j_!$ and $i^*$ are right t-exact, and $j_*$ and $i^!$ are left t-exact. If the given t-structures are bounded, so is the glued one.
\end{theorem}

Let $\cA_\cN$, $\cA$ and $\cA_\cQ$ be the three hearts. Write $\pH^0$ for the cohomology functors of the glued t-structure, and set ${}^p\!j_!:=\pH^0\circ j_!$ and ${}^p\!j_*:=\pH^0\circ j_*$ on $\cA_\cQ$. For $Y\in\cA_\cQ$ define the \emph{intermediate extension}
\[ j_{!*}Y:=\im\big({}^p\!j_!Y\to{}^p\!j_*Y\big). \]

\begin{proposition}[{\cite[\S1.4]{BBD82}}]\label{prop:glued-heart}
In the situation of \cref{thm:gluing}:
\begin{enumerate}
\item $i_*$ identifies $\cA_\cN$ with the Serre subcategory $\{X\in\cA:j^*X=0\}$ of $\cA$; this subcategory is closed under subobjects, quotients and extensions.
\item For $Y\in\cA_\cQ$ we have $j^*{}^p\!j_!Y\cong Y\cong j^*{}^p\!j_*Y\cong j^*j_{!*}Y$. The object ${}^p\!j_!Y$ has no non-zero quotient in $i_*\cA_\cN$, and ${}^p\!j_*Y$ has no non-zero subobject in $i_*\cA_\cN$. Hence $j_{!*}Y$ has neither. The kernel of ${}^p\!j_!Y\twoheadrightarrow j_{!*}Y$ and the cokernel of $j_{!*}Y\hookrightarrow {}^p\!j_*Y$ lie in $i_*\cA_\cN$.
\item The simple objects of $\cA$ are the $i_*S$ with $S$ simple in $\cA_\cN$, and the $j_{!*}S$ with $S$ simple in $\cA_\cQ$.
\item If $\cA_\cN$ and $\cA_\cQ$ are length categories, so is $\cA$.
\end{enumerate}
\end{proposition}

\begin{proof}[Comments on the proof]
(1) and (3) are proved in \cite[\S1.4]{BBD82} (see in particular 1.4.22--1.4.26 there, where $j_{!*}$ is introduced and the simple objects of $\cA$ are determined).

For (2): $j^*$ is t-exact, so $j^*{}^p\!j_!Y=H^0(j^*j_!Y)=Y$, and similarly for $j_*$. Moreover $\Hom_\cA({}^p\!j_!Y,i_*Z)=\Hom(j_!Y,i_*Z)=\Hom(Y,j^*i_*Z)=0$, and dually for ${}^p\!j_*Y$. The map ${}^p\!j_!Y\to{}^p\!j_*Y$ becomes an isomorphism after applying the exact functor $j^*|_\cA$, so its kernel and cokernel lie in $i_*\cA_\cN$ by~(1).

For (4): every $X\in\cA$ sits in an exact sequence
\[0\to i_*H^{-1}(i^*X)\to{}^p\!j_!j^*X\to X\to i_*H^0(i^*X)\to0,\]
obtained from the long exact cohomology sequence of $j_!j^*X\to X\to i_*i^*X$. The functor ${}^p\!j_!$ is right exact. By~(2), ${}^p\!j_!$ of a simple object has finite length, hence so does ${}^p\!j_!$ of any object of finite length. Thus $X$ has finite length.
\end{proof}

\begin{definition}\label{def:adm-filtration}
An \emph{admissible filtration} of $\cD$ is a chain $0=\cD_0\subset\cD_1\subset\dots\subset\cD_m=\cD$ of thick subcategories such that $\cD_{s-1}$ is admissible in $\cD_s$ for $1\le s\le m$.
\begin{itemize}
\item Its \emph{strata} are $\cT_s:=\cD_s/\cD_{s-1}$. We write $i_s,j_s$ for the functors of the recollement $(\cD_{s-1},\cD_s,\cT_s)$.
\item By abuse of notation we regard $j_{s!}$ and $j_{s*}$ as functors $\cT_s\to\cD$, by composing with the inclusion $\cD_s\subset\cD$.
\item Given t-structures on all $\cT_s$, the \emph{glued t-structure} on $\cD$ is defined inductively: the glued t-structure on $\cD_s$ is obtained from that on $\cD_{s-1}$ and the one on $\cT_s$ by \cref{thm:gluing}. (For $s=1$ we have $\cD_1=\cT_1$.)
\end{itemize}
By \cref{thm:gluing}, the inclusions $\cD_{s}\subset\cD$ are t-exact. Hence $\cA_s:=\cA\cap\cD_s$ is the heart of $\cD_s$, and it is a Serre subcategory of $\cA$.
\end{definition}

\subsection{Semi-orthogonal decompositions}
The following standard lemma is used throughout.

\begin{lemma}\label{lem:sod}
Let $\cC=\thick{\cU\cup\cV}$ for thick subcategories $\cU,\cV\subset\cC$ with $\Hom^\bullet(\cV,\cU)=0$.
\begin{enumerate}
\item Every $X\in\cC$ fits into a triangle $V_X\to X\to U_X\to V_X[1]$ with $V_X\in\cV$ and $U_X\in\cU$.
\item $X\mapsto V_X$ is right adjoint to the inclusion $\cV\subset\cC$, and $X\mapsto U_X$ is left adjoint to the inclusion $\cU\subset\cC$. Moreover $\cV^\perp\cap\cC=\cU$ and ${}^\perp\cU\cap\cC=\cV$.
\item If $\Hom^\bullet(\cV,X)=0$ then $X\in\cU$, and if $\Hom^\bullet(X,\cU)=0$ then $X\in\cV$.
\end{enumerate}
\end{lemma}

\begin{proof}
Let $\cC'\subset\cC$ be the full subcategory of objects admitting a triangle as in (1).

For $X\in\cC'$ and $V\in\cV$, the vanishing of $\Hom^\bullet(V,U_X)$ gives $\Hom(V,V_X)\cong\Hom(V,X)$. Dually, $\Hom(U_X,U)\cong\Hom(X,U)$ for $U\in\cU$. Hence the triangle is unique up to unique isomorphism, and it is functorial in $X\in\cC'$.

The subcategory $\cC'$ contains $\cU$ and $\cV$ and is closed under shifts. It is closed under cones by the $3\times3$ lemma \cite[Prop.~1.1.11]{BBD82}, applied to the morphism of triangles induced by a morphism in $\cC'$. It is closed under direct summands, because idempotents act on the functorial triangle and split in the thick subcategories $\cU$ and $\cV$. Hence $\cC'=\cC$, which proves (1); the adjunctions in (2) were shown above.

For (3): if $\Hom^\bullet(\cV,X)=0$, then $V_X\to X$ vanishes, so $U_X\cong X\oplus V_X[1]$. Then $V_X[1]$ lies in $\cU\cap\cV$, which is $0$ because $\Hom(\cV,\cU)=0$. Hence $X\cong U_X\in\cU$. The second statement is dual.

The orthogonality statements in (2) follow from (3).
\end{proof}

\subsection{Exceptional sequences}

\begin{definition}[{\cite{Bon89}; conventions of \cite{BB26}}]
An object $E$ is \emph{exceptional} if $\End^\bullet(E)=\kk$. A sequence $(E_1,\dots,E_n)$ of exceptional objects is an \emph{exceptional sequence} if $\Hom^\bullet(E_j,E_i)=0$ for $j>i$. It is \emph{full} if $\thick{E_1,\dots,E_n}=\cD$.

For an exceptional pair $(E,F)$, the left mutation $L_EF$ is defined by the triangle $\Hom^\bullet(E,F)\otimes E\to F\to L_EF$. The \emph{left dual} of $(E_1,\dots,E_n)$ is $(F_n,\dots,F_1)$, where $F_1=E_1$ and $F_i=L_{E_1}\cdots L_{E_{i-1}}E_i$.
\end{definition}

A full exceptional sequence gives the admissible filtration $\cD_s=\thick{E_1,\dots,E_s}$, with all strata equivalent to $\Db(\kk)$. In the notation of \cref{def:adm-filtration} one has $E_s=j_{s!}\kk$ and $F_s=j_{s*}\kk$ \cite[\S2.4]{BB26}. If the sequence is full, the left dual is characterised by the \emph{Hom-duality}
\begin{equation}\label{eq:hom-duality-exc}
\Hom^\bullet(E_i,F_j)\cong\delta_{ij}\kk\quad\text{(concentrated in degree $0$)}.
\end{equation}
The \emph{t-structure glued along $(E_1,\dots,E_n)$} is the one obtained by gluing the standard t-structures on $\Db(\kk)$.

\begin{theorem}[{\cite[Lemma~2.4]{BB26}}]\label{thm:BB-2.4}
Let $(E_1,\dots,E_n)$ be a full exceptional sequence with left dual $(F_n,\dots,F_1)$. The t-structure glued along it is
\[\cD^{\le0}=\{X:\Hom(X,F_s[l])=0\ \forall s,\ \forall l<0\},\qquad \cD^{\ge0}=\{X:\Hom(E_s,X[l])=0\ \forall s,\ \forall l<0\}.\]
\end{theorem}

\subsection{Highest weight categories and the Bodzenta--Bondal theorems}

\begin{definition}[{\cite{CPS88}; cf.\ \cite[\S2.1]{BB26}}]\label{def:hw}
A $\kk$-linear abelian category is \emph{Deligne finite} if it is a Hom-finite length category with finitely many simple objects and a projective generator; equivalently, it is equivalent to $\modd\Lambda$ for a finite-dimensional algebra $\Lambda$.

Let $\cA$ be Deligne finite, with simple objects $L(1),\dots,L(m)$ indexed by the totally ordered set $\{1<\dots<m\}$, projective covers $P(i)$ and injective hulls $I(i)$. Then $\cA$ is \emph{highest weight} if there are objects $\Delta(i)$ such that:
\begin{itemize}
\item[(st1)] there is an epimorphism $\Delta(i)\twoheadrightarrow L(i)$ whose kernel has composition factors $L(j)$ with $j<i$ only;
\item[(st2)] there is an epimorphism $P(i)\twoheadrightarrow\Delta(i)$ whose kernel lies in $\cF(\Delta(j):j>i)$.
\end{itemize}
The \emph{costandard objects} $\nabla(i)$ are defined dually, using injective hulls.
\end{definition}

By \cite{CPS88}, highest weight categories are exactly the module categories of quasi-hereditary algebras. \cite{BB26} also allows partial orders. Since every refinement of the partial order gives an equivalent highest weight structure \cite[Rem.~2.3]{BB26}, we restrict to total orders throughout.

\begin{definition}[{\cite{DR92}; \cite[\S2.2]{BB26}}]\label{def:univ-ext}
Let $Q,T$ be objects of an abelian category with $V:=\Ext^1(Q,T)$ finite-dimensional. The \emph{universal extension} of $Q$ by $T$ is the extension
\[0\to T\otimes V^\vee\to R\to Q\to0\]
classified by the canonical element of $\Ext^1(Q,T\otimes V^\vee)\cong V\otimes V^\vee\cong\End(V)$.

The \emph{iterated universal extension} $G=\bigoplus_{s=1}^mG_s$ of a sequence $(T_1,\dots,T_m)$ is defined recursively. For $m=1$ put $G_1=T_1$. For $m>1$, let $G'=\bigoplus_{s<m}G'_s$ be the iterated universal extension of $(T_1,\dots,T_{m-1})$. Then $G_s$ is the universal extension of $G'_s$ by $T_m$ for $s<m$, and $G_m:=T_m$.

A sequence $(E_1,\dots,E_m)$ of objects of $\cA$ is \emph{standardizable} if all $\Hom(E_i,E_j)$ and $\Ext^1(E_i,E_j)$ are finite-dimensional, and $\rad(E_i,E_j)=0=\Ext^1(E_i,E_j)$ for $i\ge j$.
\end{definition}

\begin{definition}[Universal extensions in a triangulated category]\label{def:univ-ext-tri}
Let $Q,T\in\cD$ with $V:=\Hom(Q,T[1])$ finite-dimensional. The \emph{universal extension} of $Q$ by $T$ is the object $R$ in a triangle
\[T\otimes V^\vee\to R\to Q\xrightarrow{\ \varepsilon\ }T\otimes V^\vee[1],\]
where $\varepsilon\in\Hom(Q,T\otimes V^\vee[1])\cong V\otimes V^\vee\cong\End(V)$ is the identity. It is unique up to isomorphism.

The \emph{iterated universal extension} $G=\bigoplus_sG_s$ of a sequence $(T_1,\dots,T_m)$ of objects of $\cD$ is defined recursively as in \cref{def:univ-ext}.

If $Q$ and $T$ lie in the heart of a t-structure, then so does $R$, and $R$ is the universal extension of $Q$ by $T$ in the heart.
\end{definition}

\begin{theorem}[{Dlab--Ringel \cite[Thm.~2]{DR92}; \cite[Thm.~2.1]{BB26}}]\label{thm:DR}
Let $(E_1,\dots,E_m)$ be a standardizable sequence in a $\kk$-linear abelian category $\cA$, with iterated universal extension $G$. Then $\widetilde\cA=\modd\End_\cA(G)$, with the order $1<\dots<m$, is a highest weight category with standard objects $\Hom_\cA(G,E_i)$. Moreover $\cF(E_1,\dots,E_m)\subset\cA$ and $\cF(\Delta(1),\dots,\Delta(m))\subset\widetilde\cA$ are equivalent.
\end{theorem}

\begin{theorem}[{\cite[Thm.~3.1]{BB26}}]\label{thm:BB-3.1}
Let $(E_1,\dots,E_n)$ be a full exceptional sequence in $\cD$, let $(F_n,\dots,F_1)$ be its left dual, and let $\cA$ be the heart of the t-structure glued along $(E_1,\dots,E_n)$.
\begin{enumerate}
\item $\cA$ is highest weight with standard objects $\{E_i\}$, costandard objects $\{F_i\}$ and the order induced by the sequence if and only if
\[\Hom(\textstyle\bigoplus_iE_i,\bigoplus_iE_i[l])=0\quad\text{and}\quad\Hom(\bigoplus_iF_i,\bigoplus_iF_i[l])=0\qquad\text{for all }l<0.\]
\item If $\cD$ is idempotent complete and dg enhanced and the conditions of \textup{(1)} hold, then $\cD\simeq\Db(\cA)$.
\end{enumerate}
\end{theorem}

\begin{theorem}[{\cite[Thm.~3.4]{BB26}}]\label{thm:BB-3.4}
Let $\cD$ carry a t-structure with heart $\cA$, and let $(E_1,\dots,E_n)$ be an exceptional sequence in $\cD$ with left dual $(F_n,\dots,F_1)$.
\begin{enumerate}
\item If $E_i\in\cD^{\le0}$ and $F_i\in\cD^{\ge0}$ for all $i$, then the t-structure of $\cD$ restricts to the t-structure on $\cC:=\thick{E_1,\dots,E_n}$ glued along $(E_1,\dots,E_n)$. Its heart is $\cA\cap\cC$.
\item If $E_i,F_i\in\cA$ for all $i$, then $\cA\cap\cC$ is highest weight with standard objects $\{E_i\}$ and costandard objects $\{F_i\}$.
\end{enumerate}
\end{theorem}

\begin{theorem}[{\cite[Thm.~3.8]{BB26}}]\label{thm:BB-3.8}
Let $\cA$ be an abelian category. Equivalence classes of highest weight structures on $\cA$ are in bijection with equivalence classes of full exceptional sequences in $\Db(\cA)$ such that both the sequence and its left dual consist of objects of $\cA$.
\end{theorem}

\begin{example}[{\cite[Ex.~5.11]{BB26}}]\label{ex:BB-5.11}
Let $f\colon X\to Y$ be the minimal resolution of an $A_1$ surface singularity, with exceptional curve $E$. Let $\cC_f=\{F\in\Db(X):Rf_*F=0\}$, and let $\mathcal{A}_f=\cC_f\cap\mathrm{Coh}(X)$.
\begin{itemize}
\item $\mathcal{A}_f$ is equivalent to $\modd\kk$, with additive generator $\cO_E(-1)$.
\item $\cO_E(-1)$ is $2$-spherical. In particular $\Ext^2(\cO_E(-1),\cO_E(-1))\cong\kk$, so $\Db(\mathcal{A}_f)\not\simeq\cC_f$.
\end{itemize}
In the terminology of \cref{sec:P-objects}, $\cO_E(-1)$ is a $\Pn12$-object generating $\cC_f$. This example already exhibits the phenomenon of \cref{prop:heart-P}(1).
\end{example}

\begin{remark}
Kalck's algebra \cite[Prop.~5.6]{Kal17} has a full exceptional sequence of modules but is not quasi-hereditary. Bodzenta and Bondal show that its standard t-structure is nevertheless glued along this sequence, and that the dual sequence does not consist of modules \cite[\S4.1]{BB26}. The condition on the dual sequence in \cref{thm:BB-3.1,thm:BB-3.8} is therefore essential.
\end{remark}

\subsection{Stratified algebras}

\begin{definition}[{\cite{CPS96,Dla96,AHLU00}}]\label{def:strat}
Let $\cA$ be Deligne finite, with simple objects $L(1),\dots,L(m)$ indexed by $\{1<\dots<m\}$, projective covers $P(i)$ and injective hulls $I(i)$. Define the following objects.
\begin{itemize}
\item The \emph{standard object} $\Delta(i)=P(i)/U_i$, where $U_i=\sum_{j>i}\tr_{P(j)}P(i)$. It is the largest quotient of $P(i)$ all of whose composition factors are $L(j)$ with $j\le i$.
\item The \emph{proper standard object} $\Deltabar(i)$: the largest quotient of $P(i)$ such that $[\Deltabar(i):L(i)]=1$ and all other composition factors are $L(j)$ with $j<i$.
\item The \emph{costandard object} $\nabla(i)$: the largest subobject of $I(i)$ all of whose composition factors are $L(j)$ with $j\le i$.
\item The \emph{proper costandard object} $\nablabar(i)$: the largest subobject of $I(i)$ such that $[\nablabar(i):L(i)]=1$ and all other composition factors are $L(j)$ with $j<i$.
\end{itemize}
We put $B_i:=\End_\cA(\Delta(i))$. The category $\cA$ (or an algebra $\Lambda$ with $\cA\simeq\modd\Lambda$) is
\begin{itemize}
\item \emph{standardly stratified} if $U_i\in\cF(\Delta(j):j>i)$ for every $i$, i.e.\ $P(i)$ has a $\Delta$-filtration with top factor $\Delta(i)$ and remaining factors $\Delta(j)$, $j>i$;
\item \emph{properly stratified} if moreover every $P(i)$ lies in $\cF(\Deltabar)$;
\item \emph{quasi-hereditary} if it is standardly stratified and $\Delta(i)=\Deltabar(i)$ for all $i$.
\end{itemize}
\end{definition}

\begin{remark}
Several inequivalent notions of ``stratified'' are in use: some allow partial preorders, and the notion of \cite{CPS96} is homological. For total orders, the definition above is equivalent to $\Lambda_\Lambda\in\cF(\Delta)$ \cite{Dla96,AHLU00}; Mazorchuk \cite{Maz04} calls this \emph{strongly standardly stratified}. Quasi-hereditary in the above sense is the same as highest weight in the sense of \cref{def:hw}.
\end{remark}

\begin{intuition}
The local algebra $B_i$ is the ``stratum algebra'' at $i$. Quasi-heredity means $B_i=\kk$ for all $i$. In the derived picture, $\Delta(i)$ plays the role of a projective generator of the $i$-th stratum, and $\nablabar(i)$ the role of its simple object (\cref{lem:strat-basic}). For $B_i\ne\kk$ these are genuinely different objects.
\end{intuition}

\begin{lemma}\label{lem:strat-basic}
Let $\cA$ be standardly stratified.
\begin{enumerate}
\item Let $X\in\cF(\Delta(j):j>i)$, and let $Y\in\cA$ have all composition factors in $\{L(j):j\le i\}$. Then $\Ext^l_\cA(X,Y)=0$ for all $l\ge0$.
\item $\Ext^l(\Delta(j),\Delta(i))=0$ for $j>i$ and all $l\ge0$, and $\Ext^l(\Delta(i),\Delta(i))=0$ for $l\ge1$. The algebra $B_i$ is local, and $\dim B_i=[\Delta(i):L(i)]$. In particular $\cA$ is quasi-hereditary if and only if $B_i\cong\kk$ for all $i$.
\item $\Ext^l(\Delta(j),\nablabar(i))=0$ for $j>i$ and all $l\ge0$. Furthermore $\Hom(\Delta(i),\nablabar(i))\cong\kk$ and $\Ext^l(\Delta(i),\nablabar(i))=0$ for $l\ge1$, while $\Hom(\Delta(j),\nablabar(i))=0$ for $j<i$.
\item Every $\Delta(i)$ has finite projective dimension, and
\[\Kb(\proj\cA)=\thick[\Db(\cA)]{\Delta(1),\dots,\Delta(m)}.\]
\end{enumerate}
\end{lemma}

\begin{proof}
(1) It suffices to treat $X=\Delta(j)$ with $j>i$; we use descending induction on $j$.
\begin{itemize}
\item Since $\Delta(j)$ has top $L(j)$ and $Y$ has no composition factor $L(j)$, we get $\Hom(\Delta(j),Y)=0$.
\item The exact sequence $0\to U_j\to P(j)\to\Delta(j)\to0$, together with $\Hom(P(j),Y)=0$, gives $\Ext^1(\Delta(j),Y)\cong\Hom(U_j,Y)$ and $\Ext^{l+1}(\Delta(j),Y)\cong\Ext^l(U_j,Y)$ for $l\ge1$.
\item Since $U_j\in\cF(\Delta(t):t>j)$, the induction hypothesis gives $\Ext^\bullet(U_j,Y)=0$. For $j=m$ we have $U_m=0$.
\end{itemize}

(2) The first two vanishing statements follow from~(1) applied to $Y=\Delta(i)$, using the sequence $0\to U_i\to P(i)\to\Delta(i)\to0$ as in~(1). Since $\Delta(i)$ has simple top, it is indecomposable, so $B_i$ is local. Every map $P(i)\to\Delta(i)$ vanishes on $U_i$, because $\Hom(U_i,\Delta(i))=0$ by~(1). Hence $B_i=\Hom(P(i),\Delta(i))$, whose dimension is $[\Delta(i):L(i)]$. The last statement follows, since $\Delta(i)=\Deltabar(i)$ if and only if $[\Delta(i):L(i)]=1$.

(3) The first statement is~(1) with $Y=\nablabar(i)$.
\begin{itemize}
\item The image of a non-zero map $\Delta(j)\to\nablabar(i)$ is a quotient of $\Delta(j)$ and contains $\soc\nablabar(i)=L(i)$. This is impossible for $j<i$, and forces $\dim\Hom(\Delta(i),\nablabar(i))\le\dim\Hom(P(i),\nablabar(i))=[\nablabar(i):L(i)]=1$ for $j=i$.
\item The composite $\Delta(i)\twoheadrightarrow L(i)\hookrightarrow\nablabar(i)$ is non-zero, so $\Hom(\Delta(i),\nablabar(i))\cong\kk$.
\item The vanishing of $\Ext^{l}(\Delta(i),\nablabar(i))$ for $l\ge1$ follows from $\Ext^\bullet(U_i,\nablabar(i))=0$, which holds by~(1).
\end{itemize}

(4) We use descending induction: $\Delta(m)=P(m)$, and $0\to U_i\to P(i)\to\Delta(i)\to0$ with $U_i\in\cF(\Delta(j):j>i)$. This shows that $\pd\Delta(i)<\infty$ and that $P(i)\in\thick[\Db(\cA)]{\Delta(i),\dots,\Delta(m)}$.
\end{proof}

\begin{lemma}[Inductive characterisation]\label{lem:strat-induction}
Let $\cA$ be Deligne finite, with simple objects $L(1),\dots,L(m)$ and projective covers $P(1),\dots,P(m)$. Let $\cA'\subset\cA$ be the Serre subcategory of objects without composition factor $L(m)$. For $s<m$ put
\[U_s:=\tr_{P(m)}P(s)\qquad\text{and}\qquad P'(s):=P(s)/U_s,\]
so that $P'(s)$ is a projective cover of $L(s)$ in $\cA'$. Then $\cA$ is standardly stratified (for $1<\dots<m$) if and only if both of the following hold:
\begin{enumerate}
\item[(a)] for every $s<m$, the object $U_s$ is projective, i.e.\ $U_s\cong P(m)^{\oplus q_s}$ for some $q_s\ge0$;
\item[(b)] $\cA'$ is standardly stratified (for $1<\dots<m-1$).
\end{enumerate}
In this case $\Delta(m)=P(m)$. For $s<m$, the objects $\Delta(s)$, $\Deltabar(s)$ and $\nablabar(s)$ of $\cA$ coincide with those of $\cA'$.
\end{lemma}

\begin{proof}
\emph{Comparison of $\cA$ and $\cA'$.} Let $s<m$. The objects $\Delta(s)$ and $\Deltabar(s)$ are quotients of $P(s)$ with composition factors $L(t)$, $t\le s<m$. They therefore lie in $\cA'$, are quotients of $P'(s)$, and agree with the corresponding objects of $\cA'$. Similarly, $\nablabar(s)$ lies in the largest subobject of $I(s)$ belonging to $\cA'$, which is an injective hull of $L(s)$ in $\cA'$. Finally $\Delta(m)=P(m)$, since all composition factors are $\le m$.

\emph{(a) and (b) imply standard stratification.} Let $s<m$. By (b), the kernel of $P'(s)\twoheadrightarrow\Delta(s)$ lies in $\cF(\Delta(t):s<t<m)$. The kernel of the composite $P(s)\twoheadrightarrow P'(s)\twoheadrightarrow\Delta(s)$ is an extension of this kernel by $U_s\cong\Delta(m)^{\oplus q_s}$. Hence it lies in $\cF(\Delta(t):t>s)$. For $s=m$ there is nothing to show.

\emph{Standard stratification implies (a) and (b).} We first show the following claim: every $X\in\cF(\Delta)$ has a subobject $M\in\add P(m)$ with $X/M\in\cF(\Delta(t):t<m)$. We use induction on the length of a $\Delta$-filtration. Take an exact sequence $0\to X_1\to X\to\Delta(t)\to0$, and choose $M_1\subset X_1$ for $X_1$ by induction.
\begin{itemize}
\item If $t<m$, take $M=M_1$.
\item If $t=m$, the extension $0\to X_1/M_1\to X/M_1\to\Delta(m)\to0$ splits, because $\Delta(m)=P(m)$ is projective. Let $M$ be the preimage in $X$ of the summand $\Delta(m)$. Then $0\to M_1\to M\to P(m)\to0$ also splits, so $M\in\add P(m)$, and $X/M\cong X_1/M_1$.
\end{itemize}
Now let $s<m$, and let $K_s$ be the kernel of $P(s)\twoheadrightarrow\Delta(s)$, so $K_s\in\cF(\Delta(t):t>s)$. The claim gives $M\subset K_s$ with $M\in\add P(m)$ and $K_s/M\in\cF(\Delta(t):s<t<m)$.
\begin{itemize}
\item $P(s)/M$ is an extension of $\Delta(s)$ by $K_s/M$, so it lies in $\cA'$. Hence $M\supset U_s$.
\item $M$ is generated by $P(m)$, so $M\subset U_s$.
\end{itemize}
Thus $U_s=M$, which proves (a). Moreover the kernel of $P'(s)\twoheadrightarrow\Delta(s)$ is $K_s/M\in\cF(\Delta(t):s<t<m)$, which proves (b).
\end{proof}

\begin{remark}[Heredity chains]\label{rem:heredity}
Let $\cA=\modd\Lambda$, and let $e$ be a primitive idempotent with $e\Lambda=P(m)$. Then $\bigoplus_sU_s$ corresponds to the trace ideal $J=\Lambda e\Lambda$, and $\cA'\simeq\modd\Lambda/J$. Condition (a) says that $J$ is projective as a right $\Lambda$-module. If in addition $e\Lambda e=\End(P(m))=\kk$, this is the definition of a heredity ideal, and \cref{lem:strat-induction} becomes the characterisation of quasi-hereditary algebras by heredity chains \cite{CPS88,DR92}.

So standardly stratified algebras are those admitting a chain of idempotent ideals, each projective as a right module over the corresponding quotient; the local algebras $e\Lambda e$ play the role of the stratum algebras (cf.\ \cite{Dla96,CPS96}). All inductive arguments below follow this pattern: they peel off the maximal simple object, or equivalently the last stratum.
\end{remark}



% ===============================================================
\section{Local-exceptional objects and local pairs}\label{sec:local}
% ===============================================================

\subsection{Local-exceptional objects}

\begin{definition}\label{def:local-exc}
An object $E$ of $\cD$ is \emph{local-exceptional} if $\Hom(E,E[l])=0$ for all $l\ne0$ and $\End(E)$ is a local algebra. A \emph{local-exceptional sequence} is a sequence $(E_1,\dots,E_m)$ of local-exceptional objects with $\Hom^\bullet(E_t,E_s)=0$ for $t>s$. It is \emph{full} if $\thick{E_1,\dots,E_m}=\cD$.
\end{definition}

\begin{example}\label{ex:local-exc}
\leavevmode
\begin{enumerate}
\item Exceptional objects are exactly the local-exceptional objects with $\End(E)=\kk$.
\item If $B$ is a finite-dimensional local algebra, then $B$ is a local-exceptional object of $\Db(\modd B)$ and of $\Kb(\proj B)$. More generally, a module $M$ over a finite-dimensional algebra $\Lambda$ is local-exceptional in $\Db(\modd\Lambda)$ if $\End(M)$ is local and $\Ext^l(M,M)=0$ for all $l\ge1$.
\item The standard modules of a standardly stratified algebra form a full local-exceptional sequence in $\Kb(\proj\Lambda)$ (\cref{thm:C}).
\item $\Pn n0$-like objects (\cref{def:Pnk}) are exactly the local-exceptional objects with endomorphism algebra $\kk[t]/(t^{n+1})$ (\cref{sec:P-zero}).
\end{enumerate}
\end{example}

\begin{intuition}
For an exceptional object $E$ we have $\thick E\simeq\Db(\kk)$, and $E$ is simultaneously the projective object and the simple object of the heart $\modd\kk$. For a local-exceptional object these two roles separate.
\begin{itemize}
\item $E$ behaves like the free module over $B=\End(E)$ (\cref{prop:local-exc}).
\item The ``simple'' object that $E$ should be paired with is not contained in $\thick E$.
\end{itemize}
Accordingly, the strata of the filtrations considered in \cref{sec:gluing,sec:criterion} are not generated by the objects $E_s$. Each of them contains both a projective-like and a simple-like object.
\end{intuition}

\begin{lemma}\label{lem:euler}
Let $B$ be a finite-dimensional local algebra, and let $Y\in\Kb(\proj B)$. Then $\chi(Y):=\sum_l(-1)^l\dim H^l(Y)$ is divisible by $\dim B$.
\end{lemma}

\begin{proof}
Finitely generated projective modules over a local algebra are free. So $Y$ is isomorphic to a bounded complex with terms $B^{r_l}$, and $\chi(Y)=\sum_l(-1)^lr_l\dim B$.
\end{proof}

\begin{proposition}\label{prop:local-exc}
Let $\cD$ be algebraic, let $E\in\cD$ be local-exceptional, and put $B=\End(E)$.
\begin{enumerate}
\item $\thick E\simeq\Kb(\proj B)$, with $E$ corresponding to $B$.
\item $\thick E$ carries the bounded co-t-structure with co-heart $\add E$.
\item If $B\ne\kk$, there is no $X\in\thick E$ with $\dim\Hom^\bullet(E,X)=1$. In particular, the standard t-structure of $\Db(\modd B)$ does not restrict to $\Kb(\proj B)$.
\end{enumerate}
\end{proposition}

\begin{proof}
(1) The object $E$ is a tilting object of $\thick E$, so this is \cref{thm:keller}. (2) follows from \cref{thm:silting-cot}.

(3) Under the equivalence of (1), $\Hom^l(E,X)\cong H^l(Y)$, where $Y$ is the perfect complex corresponding to $X$. If $\dim\Hom^\bullet(E,X)=1$, then $\chi(Y)=\pm1$. This is impossible by \cref{lem:euler}, since $\dim B\ge2$.

For the last assertion, let $P_1\to P_0\to\kk\to0$ be a minimal projective presentation of the simple $B$-module, and let $Y=(P_1\to P_0)$ in degrees $-1,0$. If the standard t-structure restricted to $\Kb(\proj B)$, then $H^0(Y)\cong\kk$ would be perfect. But $\Hom^\bullet(B,\kk)\cong\kk$, contradicting what we just proved.
\end{proof}

\begin{remark}\label{rem:local-exc-no-t}
For $B=\kk[t]/(t^{n+1})$ with $n\ge1$, the category $\thick E$ has no non-trivial t-structure at all. This is due to Bonfert \cite[Cor.~3.10]{Bonf26} (\cref{thm:Bonfert}(3)), and for $n=1$ to Holm, J\o rgensen and Yang (\cref{thm:HJY}). Part~(3) above is a weak form of this, valid for every local algebra $B$. It is the form needed in \cref{thm:C}.
\end{remark}


\subsection{Local pairs}
We now describe, in terms of two objects, the categories that will occur as strata of our filtrations. Throughout this subsection $\cT$ is a triangulated category (Hom-finite and idempotent complete, as always).

The following lemma describes the aisles of a bounded t-structure with a local heart in terms of two ``test objects''. It is the local ingredient in the description of the t-structure of a stratifying triple (\cref{prop:triple-t}).

\begin{lemma}\label{lem:aisles}
Let $\cT$ carry a bounded t-structure whose heart $\cH$ is a length category with a unique simple object $S$ (up to isomorphism). Let $Q\in\cH$ satisfy $\Hom(Q,M)\ne0$ for every non-zero $M\in\cH$. Then
\[\cT^{\le0}=\{X:\Hom(X,S[l])=0\ \forall l<0\},\qquad\cT^{\ge0}=\{X:\Hom(Q,X[l])=0\ \forall l<0\}.\]
Moreover $\cT^{\le0}=\Filt\{S[l]:l\ge0\}$ and $\cT^{\ge0}=\Filt\{S[l]:l\le0\}$. The hypothesis on $Q$ holds for $Q=S$, and for a projective cover of $S$ if one exists.
\end{lemma}

\begin{proof}
\textbf{Description of $\cT^{\ge0}$.} If $X\in\cT^{\ge0}$ and $l<0$, then $X[l]\in\cT^{\ge1}$, so $\Hom(Q,X[l])=0$. Conversely, suppose $X\notin\cT^{\ge0}$. By boundedness there is a minimal $d<0$ with $H^d(X)\ne0$, and we have a triangle $H^d(X)[-d]\to X\to\tau^{\ge d+1}X\to$. Shifting by $d$ and applying $\Hom(Q,-)$ gives an exact sequence
\[\Hom(Q,(\tau^{\ge d+1}X)[d-1])\to\Hom(Q,H^d(X))\to\Hom(Q,X[d]).\]
The first term vanishes because $(\tau^{\ge d+1}X)[d-1]\in\cT^{\ge2}$. The middle term is non-zero by hypothesis. Hence $\Hom(Q,X[d])\ne0$ with $d<0$.

\textbf{Description of $\cT^{\le0}$.} If $X\in\cT^{\le0}$ and $l<0$, then $S[l]\in\cT^{\ge1}$, so $\Hom(X,S[l])=0$. Conversely, suppose $X\notin\cT^{\le0}$, and let $d>0$ be maximal with $H^d(X)\ne0$. Since $X\in\cT^{\le d}$, for every $M\in\cH$ we have
\[\Hom(X,M[-d])\cong\Hom(\tau^{\ge d}X,M[-d])\cong\Hom(H^d(X),M).\]
Taking $M=S$ and using an epimorphism $H^d(X)\twoheadrightarrow S$, we obtain $\Hom(X,S[-d])\ne0$ with $-d<0$.

\textbf{The remaining statements.} The descriptions via $\Filt$ hold because $\cH=\Filt(S)$ and the t-structure is bounded. Every non-zero $M\in\cH$ has a simple subobject (isomorphic to $S$) and a simple quotient; the projection from a projective cover of $S$ onto the simple quotient lifts to $M$. This gives the last assertion.
\end{proof}

\begin{lemma}\label{lem:recognise-proj}
Let $\cT$ carry a bounded t-structure whose heart $\cH$ is a length category with a unique simple object $S$. Let $E\in\cT$ satisfy:
\begin{itemize}
\item $\Hom(E,E[l])=0$ for all $l<0$;
\item $\Hom(E,S)\cong\kk$;
\item $\Hom(E,S[l])=0$ for $l<0$ and for $l=1$.
\end{itemize}
Then $E$ lies in $\cH$ and is a projective cover of $S$.
\end{lemma}

\begin{proof}
By \cref{lem:aisles} (with $Q=S$), $E\in\cT^{\le0}$. Since $\Hom(E,S)\cong\Hom(H^0(E),S)\neq0$, we have $H^0(E)\ne0$.

Suppose $E\notin\cT^{\ge0}$, and let $a<0$ be minimal with $H^a(E)\ne0$. We use the argument of \cite[Lemma~3.12]{Bonf26}. Shifting the truncation triangle of $E$ gives a triangle $(\tau^{\ge a+1}E)[a-1]\to H^a(E)\to E[a]\to$, and hence an exact sequence
\[\Hom(E,(\tau^{\ge a+1}E)[a-1])\to\Hom(E,H^a(E))\to\Hom(E,E[a]).\]
The first term vanishes, because $(\tau^{\ge a+1}E)[a-1]\in\cT^{\ge2}$ and $E\in\cT^{\le0}$. Since $E\in\cT^{\le0}$, the middle term is
\[\Hom(E,H^a(E))\cong\Hom(H^0(E),H^a(E)).\]
This is non-zero: $H^0(E)$ and $H^a(E)$ are non-zero objects of a length category with unique simple object $S$, so there is a non-zero composite $H^0(E)\twoheadrightarrow S\hookrightarrow H^a(E)$. Hence $\Hom(E,E[a])\ne0$ with $a<0$, a contradiction.

So $E$ lies in $\cH$. Since $\Ext^1(E,S)=\Hom(E,S[1])=0$ and $S$ is the only simple object, $E$ is projective. Since $\Hom(E,S)\cong\kk$, its top is $S$.
\end{proof}

\begin{definition}[Local pairs]\label{def:local-pair}
A \emph{local pair} in $\cT$ is a pair $(\Pi,S)$ of objects of $\cT$ such that:
\begin{itemize}
\item[\textup{(LP1)}] $S$ is simple-minded and generates $\cT$: $\Hom(S,S[l])=0$ for $l<0$, $\End(S)=\kk$, and $\thick[\cT]{S}=\cT$;
\item[\textup{(LP2)}] $\Hom(\Pi,\Pi[l])=0$ for $l<0$; $\Hom(\Pi,S)\cong\kk$; and $\Hom(\Pi,S[l])=0$ for $l<0$ and for $l=1$.
\end{itemize}
We put $B:=\End(\Pi)$. The pair has \emph{Hom-duality up to degree $r$} ($1\le r\le\infty$) if $\Hom(\Pi,S[l])=0$ for $1\le l\le r$. It has \emph{full Hom-duality} if $r=\infty$, i.e.\ $\Hom^\bullet(\Pi,S)\cong\kk$, concentrated in degree $0$.
\end{definition}

\begin{proposition}\label{prop:local-pair}
Let $(\Pi,S)$ be a local pair in $\cT$.
\begin{enumerate}
\item There is a unique bounded t-structure $t_S$ on $\cT$ whose heart is $\cH:=\Filt_\cT(S)$. This heart is a length category with unique simple object $S$, and
\[\cT^{\le0}=\{X:\Hom(X,S[l])=0\ \forall l<0\},\qquad \cT^{\ge0}=\{X:\Hom(Q,X[l])=0\ \forall l<0\}\]
for $Q=S$ and for $Q=\Pi$.
\item $\Pi$ lies in $\cH$ and is a projective cover of $S$. The algebra $B=\End(\Pi)$ is local, and $\Hom(\Pi,-)$ induces an equivalence $\cH\simeq\modd B$ under which $S$ corresponds to the simple $B$-module.
\item If the pair has Hom-duality up to degree $r$, then $\Hom(\Pi,Y[l])=0$ for all $Y\in\cH$ and $1\le l\le r$.
\item Suppose the pair has full Hom-duality. Then:
\begin{itemize}
\item $\Pi$ is local-exceptional;
\item if $\cT$ is algebraic, $\real\colon\Db(\modd B)\simeq\Db(\cH)\to\cT$ is an equivalence sending $B$ to $\Pi$ and $\kk$ to $S$.
\end{itemize}
Conversely, if $\cT$ is algebraic and $\real$ is an equivalence, then the pair has full Hom-duality.
\end{enumerate}
\end{proposition}

\begin{proof}
(1) By \cref{thm:smc}, $\{S\}$ determines a bounded t-structure with heart $\Filt_\cT(S)$, which is a length category with unique simple object $S$. A bounded t-structure is determined by its heart, since its aisle consists of iterated extensions of non-negative shifts of objects of the heart; this gives uniqueness. The descriptions of the aisles follow from \cref{lem:aisles}. Note that $Q=\Pi$ is an admissible test object by~(2).

(2) This is \cref{lem:recognise-proj}. A length category with a projective generator $\Pi$ is equivalent to $\modd\End(\Pi)$ via $\Hom(\Pi,-)$.

(3) This follows by d\'evissage, since every $Y\in\cH$ is an iterated extension of copies of $S$.

(4) By~(3) with $r=\infty$, $\Hom(\Pi,\Pi[l])=0$ for $l\ge1$. Together with (LP2) and~(2), this shows that $\Pi$ is local-exceptional.

Every $X\in\cH$ is a quotient of some $\Pi^{\oplus k}$. By~(3), every morphism $X\to Y[n]$ with $X,Y\in\cH$ and $n\ge1$ is effaced by such an epimorphism. So $\real$ is an equivalence by \cref{thm:realisation}.

Conversely, if $\real$ is an equivalence, then $\Hom(\Pi,S[l])\cong\Ext^l_B(B,\kk)=0$ for $l\ge1$.
\end{proof}

\begin{example}\label{ex:local-pairs}
\leavevmode
\begin{enumerate}
\item For a finite-dimensional local algebra $B$, the pair $(B,\kk)$ in $\Db(\modd B)$ is a local pair with full Hom-duality. The case $B=\kk$ is the pair $(\kk,\kk)$ in $\Db(\kk)$, i.e.\ the stratum of an exceptional object.
\item Let $P$ be $\Pn nk$-like with $k\ge1$ and $n\ge1$, and let $\cT=\thick{P}$ (\cref{sec:P-positive}).
\begin{itemize}
\item If $k\ge2$, then $(P,P)$ is a local pair with Hom-duality up to degree $k-1$, but not up to degree $k$.
\item If $k=1$ and $n\ge2$, then $(Q,P)$ is a local pair with Hom-duality up to degree $n-1$, but not up to degree $n$; here $Q$ is the universal self-extension of $P$.
\item If $k=n=1$, no local pair $(\Pi,P)$ exists, since the heart $\Filt(P)$ has no non-zero projective objects.
\end{itemize}
\end{enumerate}
\end{example}

\begin{intuition}
In a local pair, the two roles played by $\kk$ in $\Db(\kk)$ are taken by two different objects: $\Pi$ is projective-like and $S$ is simple-like.
\begin{itemize}
\item The degrees $l\le1$ of $\Hom^\bullet(\Pi,S)$ are fixed by (LP2).
\item The degrees $l\ge2$ measure how far $\cT$ is from $\Db(\modd B)$ (\cref{prop:local-pair}(4)).
\end{itemize}
The pair is purely object-level data: the t-structure $t_S$ is determined by $S$ alone.
\end{intuition}

% ===============================================================
\section{Stratifying triples}\label{sec:triples}
% ===============================================================

In this section we introduce the object-level data underlying all our results. We show that it determines an admissible filtration, local pairs on the strata of this filtration, and a t-structure on $\cD$.

\subsection{Definition}

\begin{definition}[Stratifying triples]\label{def:triple}
Let $\cD$ be a triangulated category and $m\ge1$. A \emph{stratifying triple} $(\Delta_\bullet,\Deltabar_\bullet,N_\bullet)$ in $\cD$ consists of objects $\Delta_1,\dots,\Delta_m$, $\Deltabar_1,\dots,\Deltabar_m$ and $N_1,\dots,N_m$ of $\cD$, together with morphisms $\delta_s\colon\Deltabar_s\to N_s$. Put
\[\cD_s:=\thick{N_1,\dots,N_s}\qquad(0\le s\le m).\]
We require that the following hold for all $s,t$.
\begin{itemize}
\item[\textup{(ST1)}] \emph{Hom-duality in low degrees:} $\Hom^\bullet(\Delta_s,N_t)=0$ for $s\ne t$; $\Hom(\Delta_s,N_s)\cong\kk$; and $\Hom(\Delta_s,N_s[l])=0$ for $l<0$ and for $l=1$.
\item[\textup{(ST2)}] \emph{The simple-like sequence:} $\Hom^\bullet(N_s,N_t)=0$ for $s<t$; $\End(N_s)=\kk$; and $\Hom(N_s,N_s[l])=0$ for $l<0$.
\item[\textup{(ST3)}] \emph{Generation:} $\cD_m=\cD$.
\item[\textup{(ST4)}] \emph{Right dual:} $\cone(\delta_s)\in\cD_{s-1}$, and $\Hom^\bullet(\Deltabar_s,N_t)=0$ for $t<s$.
\item[\textup{(ST5)}] $\Hom(\Delta_s,\Delta_s[l])=0$ for $l<0$.
\end{itemize}
Further terminology:
\begin{itemize}
\item The triple has \emph{Hom-duality up to degree $r$} ($1\le r\le\infty$) if $\Hom(\Delta_s,N_s[l])=0$ for all $s$ and $1\le l\le r$. By (ST1) this always holds for $r=1$.
\item It has \emph{full Hom-duality} if $r=\infty$, i.e.\ $\Hom^\bullet(\Delta_s,N_t)\cong\delta_{st}\kk$, concentrated in degree $0$.
\item It satisfies \emph{(NH)} (no negative extensions) if
\[\Hom(\Delta_s,\Delta_t[l])=0\quad\text{and}\quad\Hom(N_s,N_t[l])=0\qquad\text{for all $s,t$ and all $l<0$.}\]
\end{itemize}
\end{definition}

\begin{intuition}
By \cref{prop:triple-filtration}, the filtration $\cD_1\subset\dots\subset\cD_m$ has strata $\cT_s$ carrying local pairs $(\Pi_s,S_s)$. The three sequences are then
\[\Delta_s=j_{s!}\Pi_s,\qquad\Deltabar_s=j_{s!}S_s,\qquad N_s=j_{s*}S_s.\]
That is, $\Delta_s$ is the projective extended by zero, $\Deltabar_s$ is the simple extended by zero, and $N_s$ is the simple extended by $j_{s*}$.

Two dualities relate the three sequences:
\begin{itemize}
\item $N_\bullet$ and $\Deltabar_\bullet$ are related by mutation;
\item $\Delta_\bullet$ and $N_\bullet$ are related by the Hom-duality (ST1), which pairs a projective object with a simple one.
\end{itemize}
Every condition in the definition only involves objects of $\cD$ and morphisms between them.
\end{intuition}

\begin{example}\label{ex:triples}
\leavevmode
\begin{enumerate}
\item A full exceptional sequence $(E_1,\dots,E_n)$ with left dual $(F_n,\dots,F_1)$ gives the stratifying triple $(E_\bullet,E_\bullet,F_\bullet)$, which has full Hom-duality (\cref{prop:exc-dual}).
\item An admissible filtration with a local pair in each stratum gives a stratifying triple, and conversely (\cref{prop:triple-filtration}).
\item Under suitable conditions, the standard, proper standard and proper costandard modules of a standardly stratified algebra form a stratifying triple with full Hom-duality (\cref{prop:dualisable}).
\item A semi-orthogonal sequence of $\mathbb{P}$-like objects with a left dual gives a stratifying triple (\cref{lem:P-triple}).
\end{enumerate}
\end{example}

\subsection{Right duals and peeling off the last stratum}

\begin{lemma}\label{lem:right-dual}
Let $N_1,\dots,N_m$ satisfy (ST2) and (ST3), and let $\delta_s\colon\Deltabar_s\to N_s$ satisfy (ST4). For every $s$:
\begin{enumerate}
\item $\Deltabar_s\in\cD_s$, $\Hom^\bullet(\Deltabar_s,\cD_{s-1})=0$, and $\cD_s=\thick{\cD_{s-1}\cup\{\Deltabar_s\}}$;
\item $\Hom^\bullet(\Deltabar_s,N_t)=0$ for $t\ne s$, and composition with $\delta_s$ induces an isomorphism $\Hom^\bullet(\Deltabar_s,N_s)\cong\End^\bullet(N_s)$;
\item the pair $(\Deltabar_s,\delta_s)$ is unique up to unique isomorphism.
\end{enumerate}
\end{lemma}

\begin{proof}
(1) Both $\cone(\delta_s)$ and $N_s$ lie in $\cD_s$, hence so does $\Deltabar_s$. Conversely $N_s\in\thick{\cD_{s-1}\cup\{\Deltabar_s\}}$, so $\cD_s=\thick{\cD_{s-1}\cup\{\Deltabar_s\}}$. The vanishing $\Hom^\bullet(\Deltabar_s,\cD_{s-1})=0$ follows from (ST4), since $\cD_{s-1}=\thick{N_1,\dots,N_{s-1}}$.

(2) For $t>s$ the vanishing holds because $\Deltabar_s\in\cD_s$ and $\Hom^\bullet(\cD_s,N_t)=0$ by (ST2). For $t=s$, apply $\Hom^\bullet(-,N_s)$ to the triangle $\Deltabar_s\to N_s\to\cone(\delta_s)\to$, and use $\Hom^\bullet(\cD_{s-1},N_s)=0$.

(3) Let $(\Deltabar'_s,\delta'_s)$ be a second choice. By (1), $\Hom^\bullet(\Deltabar_s,\cone(\delta'_s))=0$, so composition with $\delta'_s$ gives $\Hom(\Deltabar_s,\Deltabar'_s)\cong\Hom(\Deltabar_s,N_s)$. Hence $\delta_s=\delta'_s\varphi$ for a unique $\varphi$. Symmetrically $\delta'_s=\delta_s\psi$ for a unique $\psi$. By uniqueness, $\varphi\psi$ and $\psi\varphi$ are identities.
\end{proof}

\begin{lemma}[Peeling off the last stratum]\label{lem:peel}
Let $(\Delta_\bullet,\Deltabar_\bullet,N_\bullet)$ be a stratifying triple in $\cD$, and put $\cD'=\cD_{m-1}$.
\begin{enumerate}
\item $\cD'$ is admissible in $\cD$, with
\[{}^\perp\cD'=\thick{\Deltabar_m}\qquad\text{and}\qquad\cD'^\perp=\thick{N_m}.\]
More precisely, \cref{lem:sod} applied to $(\cV,\cU)=(\thick{\Deltabar_m},\cD')$ gives the left adjoint of the inclusion, and applied to $(\cV,\cU)=(\cD',\thick{N_m})$ it gives the right adjoint.
\item For $s<m$ we have $\Delta_s,\Deltabar_s\in\cD'$, and $(\Delta_{<m},\Deltabar_{<m},N_{<m})$ is a stratifying triple in $\cD'$. Moreover $\Delta_m\in{}^\perp\cD'$.
\item Let $j\colon\cD\to\cT:=\cD/\cD'$ be the quotient functor. Then $(\Pi,S):=(j^*\Delta_m,j^*N_m)$ is a local pair in $\cT$, and
\[j_!\Pi\cong\Delta_m,\qquad j_!S\cong\Deltabar_m,\qquad j_*S\cong N_m.\]
Moreover $\Hom^l_\cT(\Pi,S)\cong\Hom^l(\Delta_m,N_m)$ and $\End^\bullet_\cT(\Pi)\cong\End^\bullet(\Delta_m)$ for all $l$.
\end{enumerate}
\end{lemma}

\begin{proof}
(1) We have $\cD=\thick{\cD'\cup\{\Deltabar_m\}}$ with $\Hom^\bullet(\Deltabar_m,\cD')=0$, by \cref{lem:right-dual}(1). We also have $\cD=\thick{\cD'\cup\{N_m\}}$ with $\Hom^\bullet(\cD',N_m)=0$, by (ST2) and (ST3). Now apply \cref{lem:sod}(2) twice.

(2) Let $s<m$. By (ST1), $\Hom^\bullet(\Delta_s,N_m)=0$, so $\Delta_s\in{}^\perp\thick{N_m}=\cD'$ by \cref{lem:sod}(2). Moreover $\Deltabar_s\in\cD_s\subset\cD'$ by \cref{lem:right-dual}(1). Conditions (ST1), (ST2), (ST4) and (ST5) for indices $<m$ are inherited, and (ST3) holds for $\cD'$ by definition. Finally, $\Hom^\bullet(\Delta_m,N_t)=0$ for $t<m$ by (ST1), hence $\Delta_m\in{}^\perp\cD'$.

(3) We first identify the three objects.
\begin{itemize}
\item Since ${}^\perp\cD'$ is the essential image of $j_!$, (2) gives $\Delta_m\cong j_!\Pi$.
\item Similarly $\Deltabar_m\cong j_!j^*\Deltabar_m$, and $j^*\Deltabar_m\cong j^*N_m=S$ because $\cone(\delta_m)\in\cD'$.
\item $N_m\in\cD'^\perp$, which is the essential image of $j_*$, so $N_m\cong j_*S$.
\end{itemize}
By adjunction, $\Hom^l(\Delta_m,N_m)\cong\Hom^l(j_!\Pi,N_m)\cong\Hom^l_\cT(\Pi,j^*N_m)=\Hom^l_\cT(\Pi,S)$. Since $j_!$ and $j_*$ are fully faithful, $\End^\bullet(\Pi)\cong\End^\bullet(\Delta_m)$ and $\End^\bullet(S)\cong\End^\bullet(N_m)$.

Finally we check that $(\Pi,S)$ is a local pair. Since $\cD=\thick{\cD'\cup\{N_m\}}$, the object $S$ generates $\cT$. So (LP1) follows from (ST2), and (LP2) follows from (ST1) and (ST5).
\end{proof}

\begin{proposition}[Triples and filtrations]\label{prop:triple-filtration}
\leavevmode
\begin{enumerate}
\item Let $(\Delta_\bullet,\Deltabar_\bullet,N_\bullet)$ be a stratifying triple in $\cD$. Then $0=\cD_0\subset\cD_1\subset\dots\subset\cD_m=\cD$ is an admissible filtration (\cref{def:adm-filtration}). For each $s$, the objects $\Pi_s:=j_s^*\Delta_s$ and $S_s:=j_s^*N_s$ form a local pair in the stratum $\cT_s=\cD_s/\cD_{s-1}$, and
\[\Delta_s\cong j_{s!}\Pi_s,\qquad\Deltabar_s\cong j_{s!}S_s,\qquad N_s\cong j_{s*}S_s,\qquad\Hom^l(\Delta_s,N_s)\cong\Hom^l_{\cT_s}(\Pi_s,S_s)\quad\text{for all $l$}.\]
In particular the triple has Hom-duality up to degree $r$ if and only if all the local pairs $(\Pi_s,S_s)$ do.
\item Conversely, let $0=\cD_0\subset\dots\subset\cD_m=\cD$ be an admissible filtration with a local pair $(\Pi_s,S_s)$ in each stratum. Put $\Delta_s:=j_{s!}\Pi_s$, $\Deltabar_s:=j_{s!}S_s$ and $N_s:=j_{s*}S_s$, and let $\delta_s\colon j_{s!}S_s\to j_{s*}S_s$ be the canonical morphism. Then these form a stratifying triple, whose filtration is the given one.
\end{enumerate}
\end{proposition}

\begin{proof}
(1) Apply \cref{lem:peel} to the truncated triples $(\Delta_{\le s},\Deltabar_{\le s},N_{\le s})$ in $\cD_s$. These are stratifying triples by repeated application of \cref{lem:peel}(2).

(2) We check the five conditions.
\begin{itemize}
\item \emph{(ST1).} By adjunction, $\Hom^\bullet(j_{s!}\Pi_s,j_{t*}S_t)=0$ for $s\ne t$: for $s<t$, the object $j_{s!}\Pi_s$ lies in $\cD_{t-1}$ while $j_{t*}S_t\in\cD_{t-1}^\perp$; for $s>t$, the object $j_{t*}S_t$ lies in $\cD_{s-1}$ while $j_{s!}\Pi_s\in{}^\perp\cD_{s-1}$. For $s=t$ we get $\Hom^\bullet(j_{s!}\Pi_s,j_{s*}S_s)\cong\Hom^\bullet(\Pi_s,S_s)$. Now use (LP2).
\item \emph{(ST2) and (ST5)} follow from (LP1), (LP2) and full faithfulness of $j_{s*}$ and $j_{s!}$.
\item \emph{(ST3).} We show $\cD_s=\thick{\cD_{s-1}\cup\{j_{s*}S_s\}}$ by induction. Every $X\in\cD_s$ sits in a triangle $i_*i^!X\to X\to j_{s*}j_s^*X\to$, and $j_{s*}j_s^*X\in j_{s*}\thick[\cT_s]{S_s}=\thick{j_{s*}S_s}$.
\item \emph{(ST4).} The canonical morphism becomes an isomorphism after applying $j_s^*$. So its cone lies in $\ker j_s^*=\cD_{s-1}$. Moreover $j_{s!}S_s\in{}^\perp\cD_{s-1}$. \qedhere
\end{itemize}
\end{proof}

\begin{remark}[Consequences for the $\Delta_s$]\label{rem:Delta-semiorth}
In a stratifying triple, $\Hom^\bullet(\Delta_t,\Delta_s)=0$ for $t>s$. Indeed, $\Delta_t\in{}^\perp\cD_{t-1}$ and $\Delta_s\in\cD_s\subset\cD_{t-1}$. If the triple has full Hom-duality, then each $\Delta_s\cong j_{s!}\Pi_s$ is local-exceptional, by \cref{prop:local-pair}(4) and full faithfulness of $j_{s!}$. So in that case $(\Delta_1,\dots,\Delta_m)$ is automatically a local-exceptional sequence.
\end{remark}

\begin{remark}[Admissible subcategories and mutations]\label{rem:HD4}
Let objects $N_1,\dots,N_m$ satisfy (ST2) and (ST3). Then there exist $\Deltabar_s$ and $\delta_s$ satisfying (ST4) if and only if $\cD_{s-1}$ is left admissible in $\cD_s$ for every $s$.
\begin{itemize}
\item One direction is \cref{lem:peel}(1), applied in $\cD_s$.
\item Conversely, if $\cD_{s-1}$ is left admissible, the triangle $X'\to N_s\to i^*N_s\to$ with $X'\in{}^\perp\cD_{s-1}$ provides $\Deltabar_s:=X'$ and $\delta_s$.
\end{itemize}
In this language $\Deltabar_s$ is the right mutation of $N_s$ through $\cD_{s-1}$, which explains the term ``right dual''. Right admissibility, i.e.\ the existence of $j_{s*}$, is automatic by (ST2) and \cref{lem:sod}. So (ST4) is the only condition of admissibility type.
\end{remark}

\begin{proposition}\label{prop:exc-dual}
Let $(E_1,\dots,E_n)$ be a full exceptional sequence in $\cD$ with left dual $(F_n,\dots,F_1)$. Then $(E_\bullet,E_\bullet,F_\bullet)$ is a stratifying triple with full Hom-duality. It satisfies (NH) if and only if the conditions of \cref{thm:BB-3.1}(1) hold.
\end{proposition}

\begin{proof}
We check the conditions in turn.
\begin{itemize}
\item \emph{(ST1)} with full Hom-duality is \eqref{eq:hom-duality-exc}.
\item \emph{(ST2).} The left dual is an exceptional sequence \cite{Bon89}.
\item \emph{(ST5).} The $E_s$ are exceptional.
\item \emph{(ST3).} $F_s$ is obtained from $E_s$ by mutation through $\thick{E_1,\dots,E_{s-1}}$, via a triangle $i_*i^!E_s\to E_s\to F_s\to$. By induction on $s$ this gives $\thick{F_1,\dots,F_s}=\thick{E_1,\dots,E_s}$.
\item \emph{(ST4).} The same triangle gives $\delta_s\colon E_s\to F_s$ with cone $i_*i^!E_s[1]\in\cD_{s-1}$, and $\Hom^\bullet(E_s,F_t)=0$ for $t<s$ by \eqref{eq:hom-duality-exc}.
\end{itemize}
Finally, (NH) is the condition of \cref{thm:BB-3.1}(1).
\end{proof}

\subsection{The t-structure of a triple}
The following proposition is the only place where t-structures are glued. Everything else in the paper uses the resulting t-structure only through its description in terms of the objects of the triple.

\begin{proposition}[The t-structure of a triple]\label{prop:triple-t}
Let $(\Delta_\bullet,\Deltabar_\bullet,N_\bullet)$ be a stratifying triple in $\cD$.
\begin{enumerate}
\item The subcategories
\[\cD^{\le0}=\{X:\Hom(X,N_s[l])=0\ \forall s,\ \forall l<0\},\qquad\cD^{\ge0}=\{X:\Hom(\Delta_s,X[l])=0\ \forall s,\ \forall l<0\}\]
form a bounded t-structure on $\cD$, which we call the \emph{t-structure of the triple}. It is the t-structure glued along $(\cD_s)$ from the t-structures $t_{S_s}$ of the local pairs of \cref{prop:triple-filtration}. Moreover $\cD^{\ge0}=\{X:\Hom(\Deltabar_s,X[l])=0\ \forall s,\ \forall l<0\}$.
\item Its heart $\cA$ is a length category whose simple objects are $L_s:=j_{s!*}S_s$, for $1\le s\le m$. For each $s$, the t-structure restricts to $\cD_s$; the restriction is the t-structure of the truncated triple; and $\cA_s:=\cA\cap\cD_s$ is a Serre subcategory of $\cA$.
\item $\Delta_s,N_s\in\cA$ for all $s$ if and only if (NH) holds. Similarly, $\Deltabar_s\in\cA$ for all $s$ if and only if $\Hom(\Deltabar_s,\Deltabar_t[l])=0$ for all $s,t$ and $l<0$.
\item Suppose $\cD$ carries a t-structure $(\cD_\natural^{\le0},\cD_\natural^{\ge0})$ with $\Delta_s\in\cD_\natural^{\le0}$ and $N_s\in\cD_\natural^{\ge0}$ for all $s$. Then it is the t-structure of the triple.
\end{enumerate}
\end{proposition}

\begin{proof}
(1) By \cref{prop:triple-filtration}(1) and \cref{thm:gluing}, the glued t-structure $\tau$ exists, and it is bounded. We show by induction on $m$ that its aisles are as stated.

For $m=1$ we have $\cD=\cT_1$, and the claim is \cref{prop:local-pair}(1).

For $m\ge2$, write $i,j$ for the functors of the recollement $(\cD_{m-1},\cD,\cT_m)$. By \cref{thm:gluing}, $X\in\tau^{\le0}$ if and only if $i^*X\in\cD_{m-1}^{\le0}$ and $j^*X\in\cT_m^{\le0}$.
\begin{itemize}
\item By \cref{lem:peel}(2) and induction, $i^*X\in\cD_{m-1}^{\le0}$ means that $\Hom(i^*X,N_s[l])\cong\Hom(X,N_s[l])$ vanishes for $s<m$ and $l<0$.
\item By \cref{prop:local-pair}(1), $j^*X\in\cT_m^{\le0}$ means that $\Hom(j^*X,S_m[l])\cong\Hom(X,j_*S_m[l])=\Hom(X,N_m[l])$ vanishes for $l<0$.
\end{itemize}
Similarly, $X\in\tau^{\ge0}$ if and only if $i^!X\in\cD_{m-1}^{\ge0}$ and $j^*X\in\cT_m^{\ge0}$. By induction and \cref{prop:local-pair}(1), these say that $\Hom(\Delta_s,X[l])\cong\Hom(\Delta_s,i^!X[l])$ vanishes for $s<m$ and $l<0$, and that $\Hom(\Pi_m,j^*X[l])\cong\Hom(\Delta_m,X[l])$ vanishes for $l<0$. Using the test object $S_m$ instead of $\Pi_m$ gives the description with $\Deltabar_m=j_!S_m$.

(2) This follows from \cref{prop:glued-heart} by induction, since $i_*$ is t-exact and every stratum heart is a length category with unique simple object $S_s$.

(3) By (1), $\Delta_s\in\cD^{\ge0}$ is the first condition in (NH). Moreover $\Delta_s\in\cD^{\le0}$ holds automatically, because $\Hom^\bullet(\Delta_s,N_t)$ vanishes in negative degrees by (ST1). Dually, $N_s\in\cD^{\le0}$ is the second condition in (NH), and $N_s\in\cD^{\ge0}$ holds automatically by (ST1). For $\Deltabar_s$, use the second description of $\cD^{\ge0}$, together with \cref{lem:right-dual}(2) and the fact that $\End^\bullet(N_s)$ vanishes in negative degrees.

(4) Let $X\in\cD_\natural^{\le0}$ and $l<0$. Then $N_s[l]\in\cD_\natural^{\ge1}$, so $\Hom(X,N_s[l])=0$. Hence $\cD_\natural^{\le0}\subset\cD^{\le0}$. Dually, $\cD_\natural^{\ge0}\subset\cD^{\ge0}$. Two t-structures with both inclusions coincide: the second inclusion gives $\cD^{\le0}={}^\perp(\cD^{\ge1})\subset{}^\perp(\cD_\natural^{\ge1})=\cD_\natural^{\le0}$.
\end{proof}

% ===============================================================
\section{The induction step and Theorem A}\label{sec:gluing}
% ===============================================================

\subsection{Setup}
Throughout this section we fix the following data.
\begin{itemize}
\item $(\Delta_\bullet,\Deltabar_\bullet,N_\bullet)$ is a stratifying triple in $\cD$ satisfying (NH), and $\cA$ is the heart of its t-structure. Thus $\Delta_s,N_s\in\cA$ by \cref{prop:triple-t}(3).
\item $B_s=\End(\Delta_s)$, and $(\Pi_s,S_s)$ are the local pairs of \cref{prop:triple-filtration}.
\item $G^{(s)}=\bigoplus_{t\le s}G^{(s)}_t$ is the iterated universal extension of $(\Delta_1,\dots,\Delta_s)$ in $\cD$ (\cref{def:univ-ext-tri}). Since $\cA$ is closed under extensions, $G^{(s)}\in\cA$, and it coincides with the iterated universal extension computed in the abelian category $\cA$.
\item $r\in\{2,\infty\}$, and the triple has Hom-duality up to degree $r$.
\end{itemize}

\subsection{The induction step}

\begin{proposition}[Induction step]\label{prop:step}
Let $m\ge2$. Put $\cD'=\cD_{m-1}$ and $\cA'=\cA\cap\cD'$, and write $G'=G^{(m-1)}=\bigoplus_{s<m}G'_s$. For $s<m$ let $G_s$ be the universal extension of $G'_s$ by $\Delta_m$, so that $G^{(m)}=\bigoplus_{s<m}G_s\oplus\Delta_m$. Assume that $G'$ is a projective generator of $\cA'$ with $\Hom(G',Y[l])=0$ for all $Y\in\cA'$ and $1\le l\le r$. Then:
\begin{enumerate}
\item $\Delta_m$ is a projective cover of $L_m$ in $\cA$, and $\Hom(\Delta_m,Y[l])=0$ for all $Y\in\cA$ and $1\le l\le r$.
\item $G^{(m)}$ is a projective generator of $\cA$, and $\Hom(G^{(m)},Y[l])=0$ for all $Y\in\cA$ and $1\le l\le r$.
\item For $s<m$, the sequence $0\to\Delta_m\otimes V_s^\vee\to G_s\to G'_s\to0$ is exact in $\cA$, and $\tr_{\Delta_m}(G_s)=\Delta_m\otimes V_s^\vee$. Consequently, for every projective object $P$ of $\cA$, the trace $\tr_{\Delta_m}(P)$ is isomorphic to $\Delta_m^{\oplus q}$ for some $q$, and $P/\tr_{\Delta_m}(P)\in\add G'$ is a projective object of $\cA'$.
\item $N_m$ has simple socle $L_m$ and $[N_m:L_m]=1$. All other composition factors of $N_m$ lie in $\cA'$, and $\Ext^1_\cA(Y,N_m)=0$ for all $Y\in\cA'$.
\end{enumerate}
\end{proposition}

\begin{proof}
Write $\Delta=\Delta_m$, $N=N_m$, $L=L_m$, and let $(\Pi,S)$ be the local pair of the top stratum $\cT=\cD/\cD'$ (\cref{lem:peel}(3)). Recall that $\Ext^1_\cA=\Hom_\cD(-,-[1])$ on $\cA$, and that $\Ext^2_\cA$ embeds into $\Hom_\cD(-,-[2])$ (\cref{prop:ext-comparison}).

\textbf{Preliminaries.} We collect the facts used below.
\begin{itemize}
\item \emph{Orthogonality.} By \cref{lem:peel}(2), (ST2) and \cref{lem:peel}(3),
\begin{equation}\label{eq:step-orth}
\Hom^\bullet(\Delta,\cD')=0,\qquad\Hom^\bullet(\cD',N)=0,\qquad\Hom^l(\Delta,N)\cong\Hom^l_\cT(\Pi,S).
\end{equation}
The last group vanishes for $l<0$ and for $1\le l\le r$, and $\Hom(\Delta,N)\cong\kk$.
\item \emph{Simple objects.} By \cref{prop:triple-t}(2) and \cref{prop:glued-heart}, $\cA'$ is the Serre subcategory $\{X\in\cA:j^*X=0\}$. The simple objects of $\cA$ are $L$ and the simple objects of $\cA'$.
\item \emph{Two exact sequences.} Since $N=j_*S\in\cA$, we have ${}^p\!j_*S=N$, and \cref{prop:glued-heart}(2) gives an exact sequence
\begin{equation}\label{eq:nabla-seq}
0\to L\to N\to C\to0\qquad\text{with }C\in\cA'.
\end{equation}
Since $\Hom_\cA(\Delta,L)\cong\Hom(\Pi,j^*L)=\Hom(\Pi,S)\cong\kk$ and $L$ is simple, there is an exact sequence
\begin{equation}\label{eq:delta-seq}
0\to K\to\Delta\to L\to0\qquad\text{with }K\in\cA.
\end{equation}
\end{itemize}

\textbf{(1)} Let $1\le l\le r$. By \eqref{eq:nabla-seq} there is an exact sequence $\Hom(\Delta,C[l-1])\to\Hom(\Delta,L[l])\to\Hom(\Delta,N[l])$. Its outer terms vanish by \eqref{eq:step-orth}. Since also $\Hom(\Delta,Y[l])=0$ for $Y\in\cA'$, d\'evissage gives $\Hom(\Delta,Y[l])=0$ for all $Y\in\cA$. In particular $\Delta$ is projective in $\cA$. By \eqref{eq:step-orth} and $\Hom(\Delta,L)\cong\kk$, its top is $L$.

\textbf{(2)} For $s<m$ put $V_s:=\Hom(G'_s,\Delta[1])=\Ext^1_\cA(G'_s,\Delta)$. The universal extension is given by a triangle
\begin{equation}\label{eq:univ-tri}
\Delta\otimes V_s^\vee\to G_s\to G'_s\xrightarrow{\ \varepsilon\ }\Delta\otimes V_s^\vee[1],
\end{equation}
where $\varepsilon$ is the canonical element. We prove the vanishing for $G_s$ in five steps; the case of $\Delta$ is (1).
\begin{enumerate}[label=(\alph*)]
\item \emph{Objects of $\cD'$.} Since $\Hom^\bullet(\Delta,\cD')=0$, the triangle \eqref{eq:univ-tri} gives $\Hom(G_s,Y[l])\cong\Hom(G'_s,Y[l])$ for all $Y\in\cD'$ and all $l$. For $Y\in\cA'$ and $1\le l\le r$ this vanishes by hypothesis.
\item \emph{$\Hom(G_s,\Delta[1])=0$.} Apply $\Hom(-,\Delta[1])$ to \eqref{eq:univ-tri}. The map $\varepsilon^*\colon\Hom(\Delta\otimes V_s^\vee,\Delta)\to\Hom(G'_s,\Delta[1])=V_s$ is surjective: on $\kk\cdot\mathrm{id}\otimes V_s$ it is the identity. Moreover $\Hom(\Delta\otimes V_s^\vee,\Delta[1])=0$ by~(1).
\item \emph{Morphisms to $N$.} Since $\Hom^\bullet(G'_s,N)=0$ by \eqref{eq:step-orth}, the triangle \eqref{eq:univ-tri} gives $\Hom^\bullet(G_s,N)\cong V_s\otimes\Hom^\bullet(\Delta,N)$. This vanishes in degrees $1\le l\le r$.
\item \emph{Degrees $2\le l\le r$.} By \eqref{eq:nabla-seq} there is an exact sequence $\Hom(G_s,C[l-1])\to\Hom(G_s,L[l])\to\Hom(G_s,N[l])$. Its outer terms vanish by (a) and (c). Together with (a) and d\'evissage, $\Hom(G_s,Y[l])=0$ for all $Y\in\cA$ and $2\le l\le r$.
\item \emph{Degree $1$.} By \eqref{eq:delta-seq} there is an exact sequence $\Hom(G_s,\Delta[1])\to\Hom(G_s,L[1])\to\Hom(G_s,K[2])$. Its outer terms vanish by (b) and (d). Together with (a), $\Hom(G_s,Y[1])=0$ for all $Y\in\cA$.
\end{enumerate}
So $G^{(m)}$ is projective. It is a generator: by (a), $\Hom(G_s,L')\cong\Hom(G'_s,L')$ for $L'$ simple in $\cA'$, which is non-zero for some $s$ because $G'$ is a generator of $\cA'$; and $\Hom(\Delta,L)\ne0$. Since $\cA$ is a length category, $G^{(m)}$ is a projective generator.

\textbf{(3)} All three terms of \eqref{eq:univ-tri} lie in $\cA$, so it is a short exact sequence in $\cA$.
\begin{itemize}
\item \emph{The trace of $G_s$.} Since $\Hom(\Delta,G'_s)=0$ by \eqref{eq:step-orth}, every morphism $\Delta\to G_s$ factors through the subobject $\Delta\otimes V_s^\vee$. Hence $\tr_\Delta(G_s)\subset\Delta\otimes V_s^\vee$. Conversely, $\Delta\otimes V_s^\vee\cong\Delta^{\oplus\dim V_s}$ is generated by $\Delta$. So $\tr_\Delta(G_s)=\Delta\otimes V_s^\vee$ and $G_s/\tr_\Delta(G_s)=G'_s$.
\item \emph{An arbitrary projective $P$.} Clearly $\tr_\Delta(\Delta)=\Delta$. Since $G^{(m)}$ is a projective generator, $P\in\add G^{(m)}$. Traces are additive, and $\cA$ is Krull--Schmidt with $\Delta$ indecomposable. Hence $\tr_\Delta(P)$ is a direct summand of some $\Delta^{\oplus N}$, i.e.\ isomorphic to $\Delta^{\oplus q}$. Similarly $P/\tr_\Delta(P)\in\add G'$, which is projective in $\cA'$ by hypothesis (the case $l=1$).
\end{itemize}

\textbf{(4)} The functor $j^*$ is exact on $\cA$, kills $\cA'$ and sends $L$ to $S$. So $[N:L]$ is the length of $j^*N=S$, which is $1$. By \eqref{eq:step-orth}, $\Hom(Y,N)=0$ and $\Ext^1(Y,N)=\Hom(Y,N[1])=0$ for $Y\in\cA'$. So the socle of $N$ contains no simple object of $\cA'$, hence equals $L$.
\end{proof}

\begin{remark}[Formal versus non-formal parts]\label{rem:step-formal}
Parts (3) and (4) of \cref{prop:step} are formal consequences of the orthogonality relations \eqref{eq:step-orth} and of the universal-extension triangles. The trace condition~(3) needs no Ext-vanishing beyond the projectivity of $G'$ in $\cA'$.

The non-formal input is Ext-projectivity in~(2), and specifically degree $2$. There, Hom-duality in degree $2$ enters through step~(c), and the inductive hypothesis in degree $2$ enters through step~(d), which is needed for degree $1$ in step~(e). In terms of objects, (2) says that $G^{(m)}$ is ``tilting up to degree $r$'': $\Hom(G^{(m)},G^{(m)}[l])=0$ for $l<0$ and for $1\le l\le r$.
\end{remark}

\subsection{Theorem A}

\begin{theorem}[Theorem A]\label{thm:A}
Let $(\Delta_\bullet,\Deltabar_\bullet,N_\bullet)$ be a stratifying triple in a triangulated category $\cD$, with Hom-duality up to degree $2$. Let $\cA$ be the heart of its t-structure, let $G$ be the iterated universal extension of $(\Delta_1,\dots,\Delta_m)$ in $\cD$, and put $B_s=\End(\Delta_s)$.
\begin{enumerate}
\item $\Delta_s,N_s\in\cA$ for all $s$ if and only if (NH) holds.
\item Assume (NH). Then:
\begin{itemize}
\item $G\in\cA$ is a projective generator, with $\Hom(G,Y[l])=0$ for all $Y\in\cA$ and $l\in\{1,2\}$;
\item hence $\Hom(G,-)$ is an equivalence $\cA\simeq\modd\Lambda$, where $\Lambda:=\End_\cD(G)$;
\item $\cA$ is standardly stratified with respect to $1<\dots<m$, with standard objects $\Delta(s)=\Delta_s$, proper costandard objects $\nablabar(s)=N_s$, and $\End(\Delta(s))\cong B_s$;
\item $\cA$ is quasi-hereditary if and only if $B_s\cong\kk$ for all $s$.
\end{itemize}
\item Assume (NH), and that $\cD$ is algebraic. Then $\real\colon\Db(\cA)\to\cD$ is an equivalence if and only if the triple has full Hom-duality. In this case:
\begin{itemize}
\item $G$ is a tilting object of $\thick{\Delta_1,\dots,\Delta_m}$;
\item there is a triangle equivalence $\cD\simeq\Db(\modd\Lambda)$ sending $G\mapsto\Lambda$, $\Delta_s\mapsto\Delta(s)$ and $N_s\mapsto\nablabar(s)$, which restricts to $\thick{\Delta_1,\dots,\Delta_m}\simeq\Kb(\proj\Lambda)$;
\item $\thick{\Delta_1,\dots,\Delta_m}=\cD$ if and only if $\Lambda$ is quasi-hereditary.
\end{itemize}
\end{enumerate}
\end{theorem}

\begin{proof}
Part (1) is \cref{prop:triple-t}(3). Assume (NH). Put $r=2$, and $r=\infty$ when proving sufficiency in (3). We prove by induction on $m$ the following statement.
\begin{quote}
$(\mathrm I_m)$: $\cA$ is standardly stratified with respect to $1<\dots<m$, with $\Delta(s)\cong\Delta_s$, $\nablabar(s)\cong N_s$ and $\End(\Delta(s))\cong B_s$. Moreover, $G=G^{(m)}$ is a projective generator of $\cA$ with $\Hom(G,Y[l])=0$ for all $Y\in\cA$ and $1\le l\le r$.
\end{quote}

\emph{Base case $m=1$.} Here $\cD=\cT_1$, $(\Delta_1,N_1)$ is a local pair, and the t-structure is $t_{N_1}$ (\cref{prop:triple-t}(1)). By \cref{prop:local-pair}(2), $\cA\simeq\modd B_1$ with projective generator $G=\Delta_1$. By \cref{prop:local-pair}(3), $\Hom(\Delta_1,Y[l])=0$ for $1\le l\le r$. A category with a single simple object is standardly stratified, with $\Delta(1)=\Delta_1$ and $\nablabar(1)=\soc I(1)=N_1$.

\emph{Inductive step $m\ge2$.} By \cref{lem:peel}(2), the truncated triple in $\cD'=\cD_{m-1}$ is a stratifying triple. It satisfies (NH) and has Hom-duality up to degree $r$. By \cref{prop:triple-t}(2), its t-structure is the restriction, with heart $\cA'=\cA\cap\cD'$. So $(\mathrm I_{m-1})$ holds for it, and \cref{prop:step} applies with $G'=G^{(m-1)}$.
\begin{itemize}
\item By \cref{prop:step}(2), $G^{(m)}$ is a projective generator with the required vanishing. In particular $\cA$ is Deligne finite, and $\Hom(G,-)\colon\cA\to\modd\Lambda$ is an equivalence.
\item By \cref{prop:step}(1), $P(m)=\Delta_m$ is the projective cover of $L_m$. By \cref{prop:step}(3), condition (a) of \cref{lem:strat-induction} holds. Condition (b) is $(\mathrm I_{m-1})$. Hence $\cA$ is standardly stratified, with $\Delta(m)=\Delta_m$. For $s<m$, the objects $\Delta(s)$ and $\nablabar(s)$ are those of $\cA'$, namely $\Delta_s$ and $N_s$.
\item By \cref{prop:step}(4), $N_m$ embeds into $I(m)$. A subobject of $I(m)$ properly containing $N_m$, whose composition factors are $L(t)$ with $t<m$ together with a single $L(m)$, would contain a non-split extension of some $L(t)$, $t<m$, by $N_m$; this is excluded by \cref{prop:step}(4). So $\nablabar(m)=N_m$.
\end{itemize}
This proves $(\mathrm I_m)$, hence (2). The last assertion of (2) is \cref{lem:strat-basic}(2).

\emph{Sufficiency in (3).} Assume full Hom-duality. The induction with $r=\infty$ gives $\Hom(G,Y[l])=0$ for all $Y\in\cA$ and $l\ge1$. Let $X,Y\in\cA$, $n\ge2$ and $f\in\Hom(X,Y[n])$. Choose an epimorphism $G^{\oplus k}\twoheadrightarrow X$; its composite with $f$ vanishes. By \cref{thm:realisation}, $\real$ is fully faithful, hence an equivalence. Moreover:
\begin{itemize}
\item $\Hom(G,G[l])=0$ for $l\ne0$, and by induction (using \eqref{eq:univ-tri}) $\thick G=\thick{\Delta_1,\dots,\Delta_m}$. So $G$ is tilting in this subcategory.
\item The equivalence $\cD\simeq\Db(\cA)\simeq\Db(\modd\Lambda)$ is the composite of $\real^{-1}$ with the equivalence induced by $\Hom(G,-)\colon\cA\simeq\modd\Lambda$. It sends $G$ to $\Lambda$, and $\thick G$ to $\Kb(\proj\Lambda)$.
\item $\Kb(\proj\Lambda)=\Db(\modd\Lambda)$ if and only if $\gldim\Lambda<\infty$, i.e., by \cref{thm:C}(4), if and only if $\Lambda$ is quasi-hereditary.
\end{itemize}

\emph{Necessity in (3).} If $\real$ is an equivalence then, by \cref{lem:peel}(3), (2) and \cref{lem:strat-basic}(3), for every $s$
\[\Hom^l(\Delta_s,N_s)\cong\Ext^l_\cA(\Delta(s),\nablabar(s))=0\qquad\text{for }l\ge1.\]
Together with (ST1), this is full Hom-duality.
\end{proof}

\begin{remark}[The induction in terms of algebras]
Unwinding the induction, the chain $\cD_1\subset\dots\subset\cD_m$ gives a chain of Serre subcategories $\cA_1\subset\dots\subset\cA_m=\cA$. Under $\cA\simeq\modd\Lambda$ this corresponds to a chain of idempotent ideals of $\Lambda$. By \cref{prop:step}(3), each of these ideals is projective as a right module over the corresponding quotient algebra, i.e.\ it is a ``standardly stratifying'' ideal in the sense of \cref{rem:heredity}. The stratum algebras are the $B_s=\End(\Delta_s)$. For $B_s=\kk$ this is exactly a heredity chain.
\end{remark}

\begin{remark}[On Hom-duality in degree $2$]\label{rem:star}
In \cref{prop:step}, Hom-duality in degree $2$ for the top stratum enters only through step~(c), via $V_s\otimes\Hom(\Delta_m,N_m[2])$. It can therefore be dropped if $V_s=\Ext^1_\cA(G'_s,\Delta_m)=0$ for all $s$. For the earlier strata, it enters through the inductive hypothesis in step~(d), and hence in step~(e). \Cref{ex:k2} shows that the conclusion of \cref{thm:A}(2) may hold although Hom-duality fails in degree $2$.
\end{remark}

\begin{remark}[Filtration version]\label{rem:A-filtration}
By \cref{prop:triple-filtration}, \cref{thm:A} can equivalently be stated for an admissible filtration of $\cD$ together with local pairs $(\Pi_s,S_s)$ on its strata. The objects are then $\Delta_s=j_{s!}\Pi_s$ and $N_s=j_{s*}S_s$. In this form it contains \cref{thm:BB-3.1}: take all strata equal to $\Db(\kk)$, with the local pair $(\kk,\kk)$.
\end{remark}

\begin{remark}[Comparison with Bodzenta--Bondal]\label{rem:compare-BB}
For an exceptional sequence and the triple $(E_\bullet,E_\bullet,F_\bullet)$ of \cref{prop:exc-dual}, \cref{thm:A} becomes \cref{thm:BB-3.1}, with ``dg enhanced'' replaced by ``algebraic''. Indeed, standardly stratified with $B_s=\kk$ means quasi-hereditary, and for quasi-hereditary algebras $\nablabar=\nabla$.

The proof is different from that of \cite{BB26}. Bodzenta and Bondal first show that $G$ is a tilting generator of $\cD$ and then compare t-structures. When $B_s\neq\kk$, the object $G$ generates only $\thick{\Delta_1,\dots,\Delta_m}\subsetneq\cD$ (\cref{thm:C}), so this route is unavailable. Instead we proceed by induction on the number of strata, as for heredity chains.
\end{remark}

% ===============================================================
\section{The stratified Bodzenta--Bondal criterion}\label{sec:criterion}
% ===============================================================

\begin{definition}[Hom-dual triples]\label{def:hom-dual}
Let $(E_1,\dots,E_m)$ be a local-exceptional sequence in $\cD$.
\begin{itemize}
\item A \emph{Hom-dual sequence} for it is a sequence $(N_m,\dots,N_1)$ of objects such that:
\begin{itemize}
\item[\textup{(HD1)}] $\Hom^\bullet(E_s,N_t)\cong\delta_{st}\kk$, concentrated in degree $0$, for all $s,t$;
\item[\textup{(HD2)}] $\Hom^\bullet(N_s,N_t)=0$ for $s<t$; and for all $s$, $\End(N_s)=\kk$ and $\Hom(N_s,N_s[l])=0$ for $l<0$.
\end{itemize}
\item A \emph{Hom-dual triple} $(E_\bullet,\Deltabar_\bullet,N_\bullet)$ consists of $(E_1,\dots,E_m)$, a Hom-dual sequence $(N_m,\dots,N_1)$, and objects $\Deltabar_s$ with morphisms $\delta_s\colon\Deltabar_s\to N_s$, such that:
\begin{itemize}
\item[\textup{(HD3)}] $\cD=\thick{N_1,\dots,N_m}$;
\item[\textup{(HD4)}] $\cone(\delta_s)\in\thick{N_1,\dots,N_{s-1}}$, and $\Hom^\bullet(\Deltabar_s,N_t)=0$ for $t<s$.
\end{itemize}
\end{itemize}
\end{definition}

A Hom-dual triple is the same as a stratifying triple $(\Delta_\bullet,\Deltabar_\bullet,N_\bullet)$ with full Hom-duality and $\Delta_\bullet=E_\bullet$: conditions (HD1)--(HD4) are (ST1)--(ST4) with full Hom-duality, and (ST5) holds since the $E_s$ are local-exceptional. Conversely, by \cref{rem:Delta-semiorth}, in a stratifying triple with full Hom-duality the $\Delta_s$ automatically form a local-exceptional sequence. In particular, the semi-orthogonality of the $E_s$ and their local-exceptionality need not be assumed separately.

\begin{theorem}[Theorem B]\label{thm:B}
Let $\cD$ be an algebraic triangulated category, and let $(E_\bullet,\Deltabar_\bullet,N_\bullet)$ be a Hom-dual triple in $\cD$. Let $G$ be the iterated universal extension of $(E_1,\dots,E_m)$ in $\cD$, and put $\Lambda=\End_\cD(G)$. The following are equivalent:
\begin{enumerate}
\item $\Hom(E_s,E_t[l])=0$ and $\Hom(N_s,N_t[l])=0$ for all $s,t$ and all $l<0$;
\item $G$ is a tilting object of $\thick{E_1,\dots,E_m}$, the algebra $\Lambda$ is standardly stratified with respect to $1<\dots<m$, and there is a triangle equivalence $\Phi\colon\cD\simeq\Db(\modd\Lambda)$ with
\[\Phi(G)\cong\Lambda,\qquad\Phi(E_s)\cong\Delta(s),\qquad\Phi(N_s)\cong\nablabar(s).\]
\end{enumerate}
In this case:
\begin{itemize}
\item $\End(\Delta(s))\cong\End(E_s)$, and $\Lambda$ is quasi-hereditary if and only if all $E_s$ are exceptional;
\item $\Phi$ restricts to $\thick{E_1,\dots,E_m}\simeq\Kb(\proj\Lambda)$, which is a proper subcategory unless $\Lambda$ is quasi-hereditary;
\item $\Phi(\Deltabar_s)$ is the object $j_{s!}\kk$ of $\Db(\modd\Lambda)$, a ``derived proper standard object''.
\end{itemize}
\end{theorem}

\begin{proof}
\emph{(1)$\Rightarrow$(2).} Apply \cref{thm:A}(2),(3) to the stratifying triple $(E_\bullet,\Deltabar_\bullet,N_\bullet)$, which has full Hom-duality. The description of $\Phi(\Deltabar_s)$ follows from \cref{prop:triple-filtration}(1).

\emph{(2)$\Rightarrow$(1).} The objects $\Phi(E_s)$ and $\Phi(N_s)$ are modules, so they have no negative extensions.

The remaining assertions are \cref{thm:A}(2),(3).
\end{proof}

\begin{remark}[The t-structure in Theorem B]\label{rem:G-tilting}
The statement of \cref{thm:B} does not mention t-structures. In its proof, the t-structure of the triple (\cref{prop:triple-t}) corresponds under $\Phi$ to the standard t-structure of $\Db(\modd\Lambda)$. Since $\Phi(G)=\Lambda$, it is the \emph{$G$-tilting t-structure}:
\[\cD^{\le0}=\{X:\Hom(G,X[l])=0\ \forall l>0\},\qquad\cD^{\ge0}=\{X:\Hom(G,X[l])=0\ \forall l<0\}.\]
Its heart is $\{X:\Hom(G,X[l])=0\ \forall l\ne0\}\simeq\modd\Lambda$.
\end{remark}

\begin{remark}[Silting viewpoint]\label{rem:silting-view}
The object $G$ is built from the stratum projectives $\Pi_s$ by universal extensions, i.e.\ it is a silting object glued along the filtration of $\thick{E_1,\dots,E_m}$ by the $\thick{E_{\le s}}$. On this perfect part the natural adjoints are $i^*$ and $j_!$, given by the $E_s$. This is the setting of gluing silting objects \cite{LVY14}. \Cref{thm:B} can be read as follows: under (1), the glued silting object is tilting, and its Koenig--Yang dual structure on $\cD$ (\cref{thm:KY}) is described by the simple-like objects $N_s$. We have not tried to derive the theorem from silting gluing directly; see \cref{q:silting}.
\end{remark}

\begin{remark}[Comparison with Bodzenta--Bondal]\label{rem:B-vs-BB}
For a full exceptional sequence and the triple $(E_\bullet,E_\bullet,F_\bullet)$ of \cref{prop:exc-dual}, \cref{thm:B} becomes \cref{thm:BB-3.1}. The new features in the general case are the following.
\begin{enumerate}
\item The Hom-dual sequence is extra data: the left dual pairs with $E_s$ via $\End(E_s)$ rather than $\kk$.
\item The ambient category is generated by the Hom-dual sequence; the $E_s$ generate only $\Kb(\proj\Lambda)$ (\cref{thm:C}).
\item The right dual of $N_\bullet$ is a third sequence $\Deltabar_\bullet$, of ``proper standard'' objects, instead of $E_\bullet$ itself.
\end{enumerate}
\end{remark}

% ===============================================================
\section{Standardly stratified algebras}\label{sec:stratified}
% ===============================================================

\subsection{Standard modules form local-exceptional sequences}

\begin{theorem}\label{thm:C}
Let $\Lambda$ be a standardly stratified algebra with respect to $1<\dots<m$, and let $B_i=\End(\Delta(i))$.
\begin{enumerate}
\item $(\Delta(1),\dots,\Delta(m))$ is a full local-exceptional sequence in $\Kb(\proj\Lambda)$ with $\End(\Delta(i))=B_i$.
\item In $\Db(\modd\Lambda)$ we have $\Hom^\bullet(\Delta(i),\nablabar(i))\cong\kk$ (in degree $0$), and $\Hom^\bullet(\Delta(j),\nablabar(i))=0$ for $j>i$.
\item If $B_i\ne\kk$, then $\nablabar(i)\notin\Kb(\proj\Lambda)$, i.e.\ $\pd\nablabar(i)=\infty$.
\item $\gldim\Lambda<\infty$ if and only if $\Lambda$ is quasi-hereditary, if and only if $B_i\cong\kk$ for all $i$.
\end{enumerate}
\end{theorem}

\begin{proof}
(1) and (2) are \cref{lem:strat-basic}.

(3) Suppose $X:=\nablabar(i)\in\Kb(\proj\Lambda)=\thick[\Kb(\proj\Lambda)]{\Delta(1),\dots,\Delta(m)}$.
\begin{itemize}
\item Apply \cref{lem:sod}(3) in $\Kb(\proj\Lambda)$ with
\[\cU=\thick[\Kb(\proj\Lambda)]{\Delta(1),\dots,\Delta(i)},\qquad\cV=\thick[\Kb(\proj\Lambda)]{\Delta(i+1),\dots,\Delta(m)}.\]
Since $\Hom^\bullet(\cV,X)=0$ by~(2), we get $X\in\cU$.
\item Apply \cref{lem:sod}(1) to $\cU$, with $\cU'=\thick[\Kb(\proj\Lambda)]{\Delta(1),\dots,\Delta(i-1)}$ and $\cV'=\thick[\Kb(\proj\Lambda)]{\Delta(i)}$. This gives a triangle $V\to X\to U\to V[1]$ with $V\in\thick[\Kb(\proj\Lambda)]{\Delta(i)}$ and $U\in\cU'$.
\item Since $\Hom^\bullet(\Delta(i),\cU')=0$, we get $\Hom^\bullet(\Delta(i),V)\cong\Hom^\bullet(\Delta(i),X)\cong\kk$.
\end{itemize}
This contradicts \cref{prop:local-exc}(3) applied to $E=\Delta(i)$ in the algebraic category $\Db(\modd\Lambda)$.

(4) If $\gldim\Lambda<\infty$, every module is perfect, so $B_i=\kk$ for all $i$ by~(3). Then $\Lambda$ is quasi-hereditary by \cref{lem:strat-basic}(2). Conversely, quasi-hereditary algebras have finite global dimension \cite{CPS88,DR92}.
\end{proof}

\begin{remark}
Statement~(4) is presumably well known to experts (cf.\ \cite{AHLU00,Dla96}). We include the short argument because it isolates exactly the obstruction discussed in the introduction. In a genuinely stratified situation, the standard objects generate only $\Kb(\proj\Lambda)$, while their Hom-dual partners $\nablabar(i)$ lie outside it.
\end{remark}

\subsection{The converse}
We now ask which standardly stratified algebras arise from \cref{thm:B}.

Let $\cA_s\subset\modd\Lambda$ be the Serre subcategory of modules whose composition factors lie in $\{L(1),\dots,L(s)\}$. We say that $\Lambda$ satisfies the \emph{$\nablabar$-condition} if, for every $i$, there is an exact sequence
\[0\to\nablabar(i)\to I(i)\to C_i\to0\qquad\text{with } C_i\in\cF(\nablabar(t):t>i).\]

\begin{lemma}\label{lem:nablabar}
Let $\Lambda$ be standardly stratified.
\begin{enumerate}
\item $\End(\nablabar(i))=\kk$.
\item $\thick[\Db(\modd\Lambda)]{\nablabar(1),\dots,\nablabar(s)}=\Db_{\cA_s}(\modd\Lambda)$, the subcategory of complexes with cohomology in $\cA_s$. In particular $\Db(\modd\Lambda)=\thick[\Db(\modd\Lambda)]{\nablabar(1),\dots,\nablabar(m)}$.
\item If $\Lambda$ satisfies the $\nablabar$-condition, then $(\nablabar(m),\dots,\nablabar(1))$ is a Hom-dual sequence for $(\Delta(1),\dots,\Delta(m))$ in $\Db(\modd\Lambda)$ satisfying (HD3); i.e.\ (HD1)--(HD3) hold.
\end{enumerate}
\end{lemma}

\begin{proof}
(1) Let $f$ be a non-zero endomorphism of $\nablabar(i)$. If $\ker f$ contained $\soc\nablabar(i)=L(i)$, then $\im f\ne0$ would have no composition factor $L(i)$, since $[\nablabar(i):L(i)]=1$. But $\im f$ contains the socle $L(i)$, a contradiction. Hence $\ker f=0$ (the socle is essential), so $f$ is an automorphism. Thus $\End(\nablabar(i))$ is a division algebra, hence equal to $\kk$.

(2) The subcategory $\Db_{\cA_s}(\modd\Lambda)$ is $\thick[\Db(\modd\Lambda)]{L(1),\dots,L(s)}$. Moreover $\nablabar(t)$ has composition factors $L(t)$ (once) and $L(u)$ with $u<t$. By induction on $s$,
\[\thick[\Db(\modd\Lambda)]{\nablabar(1),\dots,\nablabar(s)}=\thick[\Db(\modd\Lambda)]{L(1),\dots,L(s)}.\]

(3) We first prove the dual of \cref{lem:strat-basic}(1): if $X\in\cF(\nablabar(t):t>i)$ and $Y$ has composition factors in $\{L(u):u\le i\}$, then $\Ext^\bullet(Y,X)=0$. It suffices to take $X=\nablabar(t)$ with $t>i$, and we use descending induction on $t$.
\begin{itemize}
\item $\Hom(Y,\nablabar(t))=0$, since the socle $L(t)$ is not a composition factor of $Y$.
\item $\Hom(Y,I(t))=0$ for the same reason.
\item The sequence $0\to\nablabar(t)\to I(t)\to C_t\to0$ gives $\Ext^1(Y,\nablabar(t))\cong\Hom(Y,C_t)$ and $\Ext^{l+1}(Y,\nablabar(t))\cong\Ext^l(Y,C_t)$ for $l\ge1$.
\item These groups vanish by induction, since $C_t\in\cF(\nablabar(u):u>t)$.
\end{itemize}
Now we check the conditions.
\begin{itemize}
\item (HD1) for $j\ge i$ is \cref{lem:strat-basic}(3). For $j<i$, we have $\Hom(\Delta(j),\nablabar(i))=0$ by \cref{lem:strat-basic}(3). The sequence $0\to\nablabar(i)\to I(i)\to C_i\to0$ gives $\Ext^1(\Delta(j),\nablabar(i))\cong\Hom(\Delta(j),C_i)$ and $\Ext^{l+1}(\Delta(j),\nablabar(i))\cong\Ext^l(\Delta(j),C_i)$ for $l\ge1$. These vanish by the dual statement, because $\Delta(j)$ has composition factors $L(u)$ with $u\le j<i$.
\item (HD2): for $s<t$, take $X=\nablabar(t)$ and $Y=\nablabar(s)$ in the dual statement. The remaining parts of (HD2) are (1) and the fact that modules have no negative extensions.
\item (HD3) is (2). \qedhere
\end{itemize}
\end{proof}

\begin{remark}
The $\nablabar$-condition is the injective counterpart of standard stratification. We expect that it holds for every standardly stratified algebra: it is a form of BGG reciprocity, cf.\ \cite{AHLU00,Dla96}. We have not checked this and do not use it. For quasi-hereditary algebras it is the classical fact that injective modules have $\nabla$-filtrations \cite{CPS88,DR92}.
\end{remark}

\begin{proposition}\label{prop:dualisable}
Let $\Lambda$ be standardly stratified and satisfy the $\nablabar$-condition, and suppose that
\[\Ext^l_\Lambda(\Deltabar(s),\nablabar(t))=0\qquad\text{for all }t<s\text{ and all }l.\]
Let $\delta_s$ be the composite $\Deltabar(s)\twoheadrightarrow L(s)\hookrightarrow\nablabar(s)$. Then $\big((\Delta(s))_s,(\Deltabar(s))_s,(\nablabar(s))_s\big)$, with the morphisms $\delta_s$, is a Hom-dual triple in $\Db(\modd\Lambda)$.
\end{proposition}

\begin{proof}
By \cref{lem:nablabar}(3), conditions (HD1)--(HD3) hold. For (HD4), recall that $[\Deltabar(s):L(s)]=1=[\nablabar(s):L(s)]$, and that all other composition factors of $\Deltabar(s)$ and $\nablabar(s)$ are $L(t)$ with $t<s$. So the kernel $\rad\Deltabar(s)$ and the cokernel $\nablabar(s)/L(s)$ of $\delta_s$ lie in $\cA_{s-1}$. Hence $\cone(\delta_s)$, whose cohomology consists of these two modules, lies in $\Db_{\cA_{s-1}}(\modd\Lambda)=\cD_{s-1}$, by \cref{lem:nablabar}(2). The Hom condition in (HD4) is the hypothesis.
\end{proof}

\begin{definition}
Let $\cA$ be a $\kk$-linear abelian category such that $\Db(\cA)$ is Hom-finite.
\begin{itemize}
\item A \emph{standardly stratified structure} on $\cA$ is a total order $L_1<\dots<L_m$ on the isomorphism classes of simple objects for which $\cA$ is Deligne finite and standardly stratified.
\item It is \emph{dualisable} if there are objects $\Deltabar_s\in\Db(\cA)$ and morphisms $\delta_s$ making $\big((\Delta(s))_s,(\Deltabar_s)_s,(\nablabar(s))_s\big)$ a Hom-dual triple in $\Db(\cA)$. By \cref{lem:right-dual}(3), the $\Deltabar_s$ are then unique.
\end{itemize}
\end{definition}

\begin{corollary}[Analogue of {\cref{thm:BB-3.8}}]\label{cor:BB38}
Let $\cA$ be a $\kk$-linear abelian category such that $\Db(\cA)$ is Hom-finite. The assignment
\[(L_1<\dots<L_m)\longmapsto\big((\Delta(1),\dots,\Delta(m)),\,(\nablabar(m),\dots,\nablabar(1))\big)\]
is a bijection between:
\begin{enumerate}
\item dualisable standardly stratified structures on $\cA$; and
\item isomorphism classes of pairs $(E_\bullet,N_\bullet)$ of objects of $\cA$ that extend to a Hom-dual triple $(E_\bullet,\Deltabar_\bullet,N_\bullet)$ in $\Db(\cA)$.
\end{enumerate}
In particular every dualisable standardly stratified structure is recovered by gluing. For exceptional sequences this is \cref{thm:BB-3.8}.
\end{corollary}

\begin{proof}
By \cref{thm:C}(1) and the definition, the assignment maps (1) to (2). It is injective, since the order is recovered from the sequence: $L_s$ is the top of $\Delta(s)$.

For surjectivity, let $(E_\bullet,\Deltabar_\bullet,N_\bullet)$ be a Hom-dual triple in $\Db(\cA)$ with $E_s,N_s\in\cA$.
\begin{itemize}
\item Since $E_s\in\Db(\cA)^{\le0}$ and $N_s\in\Db(\cA)^{\ge0}$ for the standard t-structure, \cref{prop:triple-t}(4) shows that the standard t-structure is the t-structure of the triple. In particular its heart is $\cA$.
\item Objects of $\cA$ have no negative extensions, so (NH) holds. Hence \cref{thm:A}(2) shows that $\cA$ is standardly stratified for the order $L_1<\dots<L_m$, where $L_s$ is the top of $E_s$, with $\Delta(s)=E_s$ and $\nablabar(s)=N_s$.
\end{itemize}
This structure is dualisable via $\Deltabar_\bullet$, and it maps to $(E_\bullet,N_\bullet)$.

For a quasi-hereditary structure, the standard objects form a full exceptional sequence in $\Db(\cA)$ whose left dual consists of the costandard objects $\nabla(s)=\nablabar(s)$ \cite[\S3]{BB26}. So the structure is dualisable with $\Deltabar_\bullet=\Delta_\bullet$, by \cref{prop:exc-dual}, and we recover \cref{thm:BB-3.8}.
\end{proof}

\begin{question}\label{q:HD4}
Is every standardly stratified structure dualisable? More concretely:
\begin{enumerate}
\item[(a)] Does every standardly stratified algebra satisfy the $\nablabar$-condition?
\item[(b)] Is $\Ext^\bullet_\Lambda(\Deltabar(s),\nablabar(t))=0$ for all $t<s$?
\end{enumerate}
By \cref{prop:dualisable}, positive answers to (a) and (b) imply dualisability, with the proper standard modules as right dual.
\end{question}

Some evidence and remarks:
\begin{itemize}
\item For quasi-hereditary algebras, $\Deltabar=\Delta$, $\nablabar=\nabla$, and both (a) and (b) are classical.
\item For properly stratified algebras we expect (b) to hold. On the top stratum, $\Deltabar(m)$ should be $j_{m!}$ of the simple $B_m$-module, and then $\Ext^\bullet(\Deltabar(m),i_*Y)=0$ by adjunction; one would then proceed by induction, peeling off the top stratum.
\item If (b) fails, the structure may still be dualisable, with $\Deltabar_s$ a complex rather than a module.
\end{itemize}

% ===============================================================
\section{Application: \texorpdfstring{$\mathbb{P}^n[k]$}{Pn[k]}-like objects}\label{sec:P-objects}
% ===============================================================

$\Pn nk$-like objects generalise exceptional objects by allowing self-extensions in the degrees $k,2k,\dots,nk$. In this section we explain how they fit into the general framework.
\begin{itemize}
\item For $k=0$ they are local-exceptional objects (\cref{sec:P-zero}), so \cref{thm:B,thm:C} apply.
\item For $k\ge1$ they are simple-minded, and they are the second member of local pairs (\cref{sec:P-positive}). Semi-orthogonal sequences of them give stratifying triples, to which \cref{thm:A} applies (\cref{sec:P-sequences}).
\end{itemize}
\Cref{tab:summary} summarises the outcome.

\begin{table}[ht]
\centering
\renewcommand{\arraystretch}{1.25}
\small
\begin{tabular}{@{}llll@{}}
\toprule
stratum, local pair & heart of the stratum & glued heart $\cA$ & $\Db(\cA)\simeq\cD$? \\
\midrule
$\Db(\kk)$, $(\kk,\kk)$ & $\modd\kk$ & highest weight \cite{BB26} & yes \\
$\Db(\modd B)$, $(B,\kk)$ & $\modd B$ & std.\ stratified (Thm.~\ref{thm:A}, \ref{thm:B}) & yes \\
$\Pn nk$, $k\ge 3$; $(P,P)$ & $\modd\kk$ & highest weight (Thm.~\ref{thm:Pseq}) & no \\
$\Pn nk$, $k=2$; $(P,P)$ & $\modd\kk$ & open (Question~\ref{q:k2}) & no$^{\ast}$ \\
$\Pn n1$, $n\ge 3$; $(Q,P)$ & $\modd\kk[x]/(x^2)$ & std.\ stratified (Thm.~\ref{thm:Pseq}) & no \\
$\Pn 21$; $(Q,P)$ & $\modd\kk[x]/(x^2)$ & open (Question~\ref{q:k2}) & no$^{\ast}$ \\
$\Pn 11$; no local pair & nilpotent $\kk[x]$-modules & not Deligne finite (Prop.~\ref{prop:P11}) & --- \\
$\Pn n0$ & no t-structure \cite{Bonf26} & --- (use Thm.~\ref{thm:B} instead) & --- \\
\bottomrule
\end{tabular}
\caption{Strata and the resulting glued heart. In the rows with $\Pn nk$ we assume $n\ge1$, and the stratum is $\thick{P}$. ``Std.\ stratified'' means standardly stratified, and ``no t-structure'' means no non-trivial t-structure. A $\Pn n0$-like object enters \cref{thm:B} as a local-exceptional object $E_s$; its stratum is generated by the Hom-dual object $N_s$ and is equivalent to $\Db(\modd\kk[t]/(t^{n+1}))$. The asterisk means ``no, whenever $\cA$ is a module category of the expected type''.}\label{tab:summary}
\end{table}


\subsection{Definitions and examples}

\begin{definition}[{\cite{Kru18}, \cite[Def.~2.1]{HK19}, \cite[Def.~2.1]{Bonf26}}]\label{def:Pnk}
Let $n\ge0$ and $k\in\mathbb{Z}$. An object $P$ of $\cD$ is \emph{$\Pn nk$-like} if there is an isomorphism of graded algebras
\[\End^\bullet(P)\cong\kk[t]/(t^{n+1}),\qquad\deg t=k.\]
It is a \emph{$\Pn nk$-object} if moreover $\Hom^\bullet(P,X)$ is finite-dimensional for all $X$ and $P$ is $nk$-Calabi--Yau, i.e.\ there are isomorphisms $\Hom(P,X)\cong\Hom(X,P[nk])^\vee$, functorial in $X\in\cD$.
\end{definition}

As in \cite{Bonf26} we allow arbitrary $k\in\mathbb{Z}$. For $k=0$ the condition says that $\End(P)\cong\kk[t]/(t^{n+1})$ and $\Hom(P,P[l])=0$ for $l\ne0$.

\begin{example}\label{ex:Pnk}
\leavevmode
\begin{enumerate}
\item The $\Pn0k$-like objects are exactly the exceptional objects.
\item $\Pn1k$-like objects are the \emph{$k$-spherelike} objects. $\Pn1k$-objects are the $k$-spherical objects of Seidel and Thomas \cite{ST01}; in the terminology of \cite{HJY13} they are \emph{$w$-spherical} with $w=k$.
\item $\Pn n2$-objects are the $\mathbb{P}^n$-objects of Huybrechts and Thomas \cite{HT06}. For example, the structure sheaf of a compact hyperk\"ahler manifold of dimension $2n$ is one.
\item Krug \cite{Kru18} introduced the notation $\Pn nk$ and gave examples of varieties whose structure sheaves are $\Pn nk$-objects; see also \cite{HK19}.
\item ($k=0$) If $B=\kk[t]/(t^{n+1})$, then $B$ is a $\Pn n0$-like object of $\Db(\modd B)$. More generally, a module $M$ with $\End(M)\cong\kk[t]/(t^{n+1})$ and $\Ext^{>0}(M,M)=0$ is $\Pn n0$-like. Standard modules of stratified algebras provide examples (\cref{thm:C}).
\end{enumerate}
\end{example}

\begin{remark}[{\cite[Rem.~2.2]{HK19}, \cite[Lemma~2.5]{Bonf26}}]
A $\Pn nk$-like object $P$ is already a $\Pn nk$-object of the subcategory $\thick P$; in fact, for $\cD$ algebraic, $\thick P$ is $nk$-Calabi--Yau. Hence for questions concerning $\thick P$ alone, the Calabi--Yau condition is irrelevant. In sequences of such objects the Calabi--Yau condition refers to the ambient category and is not preserved by mutations, so we work with ``$\Pn nk$-like'' throughout.
\end{remark}

\begin{intuition}
For an exceptional object $E$ one has $\thick E\simeq\Db(\kk)$, and $E$ is at the same time the simple object, the projective object and the injective object of the heart $\modd\kk$. A $\Pn nk$-like object is a ``thickening'' with self-extensions in degrees $k,2k,\dots,nk$, and these three roles separate:
\begin{itemize}
\item For $k\ge1$, $P$ has no negative self-extensions and $\End(P)=\kk$. It is \emph{simple-minded} and determines a t-structure.
\item For $k\le0$, $P$ has no positive self-extensions, so it is \emph{silting} and determines a co-t-structure. For $k=0$ it behaves like the free module over the local algebra $\kk[t]/(t^{n+1})$.
\end{itemize}
This is the dichotomy of \cref{thm:HJY,thm:Bonfert}; the standard t-structure and the standard weight structure on $\thick P$ are obtained in exactly this way in \cite[\S3.2]{Bonf26}.
\end{intuition}

\subsection{Known classification results}

\begin{theorem}[{Holm--J\o rgensen--Yang \cite{HJY13}}]\label{thm:HJY}
Let $w\in\mathbb{Z}$, and let $\cT$ be the $\kk$-linear algebraic triangulated category generated by a $w$-spherical object.
\begin{itemize}
\item For $w\le0$, $\cT$ has no non-trivial t-structures, but it has one family of non-trivial co-t-structures.
\item For $w\ge1$, $\cT$ has no non-trivial co-t-structures, but it has one family of non-trivial t-structures.
\end{itemize}
For $w=0$, $\cT$ is the compact derived category of the dual numbers.
\end{theorem}

Here ``non-trivial'' excludes $(\cT,0)$ and $(0,\cT)$, and a ``family'' consists of the shifts of a single structure.

\begin{theorem}[{Hochenegger--Krug \cite[Lemma~2.8]{HK19}; Bonfert \cite[Props.~2.3 and~2.4]{Bonf26}}]\label{thm:formality}
For all $n\ge0$ and $k\in\mathbb{Z}$, the graded algebra $A_{n,k}:=\kk[x]/(x^{n+1})$ with $\deg x=k$ is intrinsically formal. Consequently, let $\cD$ be an idempotent complete algebraic triangulated category and $P\in\cD$ a $\Pn nk$-like object. Then $\thick P\simeq\operatorname{perf}(A_{n,k})$, where $A_{n,k}$ is regarded as a dg algebra with zero differential, and $P$ corresponds to $A_{n,k}$.
\end{theorem}

\begin{theorem}[{Bonfert \cite[Thm.~3.3]{Bonf26}}]\label{thm:Bonfert}
Let $\cT$ be the triangulated category generated by a $\Pn nk$-object $P$.
\begin{enumerate}
\item If $n=0$, every non-trivial t-structure on $\cT$ is a shift of the standard t-structure, and every non-trivial weight structure on $\cT$ is a shift of the standard weight structure.
\item If $n>0$ and $k>0$, every non-trivial t-structure on $\cT$ is a shift of the standard t-structure, whose heart is $\Filt(P)$. There is no non-trivial weight structure on $\cT$.
\item If $n>0$ and $k\le0$, every non-trivial weight structure on $\cT$ is a shift of the standard weight structure, whose co-heart is $\add P$. There is no non-trivial t-structure on $\cT$.
\end{enumerate}
\end{theorem}

For $n=1$ this recovers \cref{thm:HJY}. By \cref{thm:formality} and the preceding remark, \cref{thm:Bonfert} applies to $\thick P$ for every $\Pn nk$-like object $P$ of an idempotent complete algebraic triangulated category.

\begin{remark}[{Koszul dual models \cite[\S4]{Bonf26}}]\label{rem:Bonfert-Koszul}
Bonfert also computes the dg Koszul dual $A_{n,k}^{!,\mathrm{dg}}=\RHom_{A_{n,k}}(\kk,\kk)$ \cite[Thm.~4.9, Props.~4.11 and~4.12]{Bonf26}:
\begin{itemize}
\item for $n=1$ and $k\ne1$ it is $\kk[s]$ with $\deg s=1-k$;
\item for $n=k=1$ it is $\kk[[s]]$ with $\deg s=0$;
\item for $n>1$ it is a non-formal dg algebra with cohomology $\kk[s,t]/(s^2)$, where $\deg s=1-k$ and $\deg t=2-(n+1)k$.
\end{itemize}
For $k=0$ this is the Ext-algebra $\Ext^\bullet_B(\kk,\kk)$ of the simple module over $B=\kk[t]/(t^{n+1})$. It is the graded endomorphism algebra of the ``simple-like'' objects $\nablabar$ that will be paired with $\Pn n0$-like objects in \cref{sec:gluing,sec:stratified}.
\end{remark}

\subsection{The case \texorpdfstring{$k=0$}{k=0}: \texorpdfstring{$\mathbb{P}^n[0]$}{Pn[0]}-like objects are local-exceptional}\label{sec:P-zero}

A $\Pn n0$-like object is the same as a local-exceptional object $E$ with $\End(E)\cong\kk[t]/(t^{n+1})$. In this case the general results specialise as follows.
\begin{itemize}
\item By \cref{thm:formality} and \cref{thm:Bonfert}(3), $\thick{E}\simeq\operatorname{perf}(\kk[t]/(t^{n+1}))$ carries, up to shift, a unique non-trivial weight structure, with co-heart $\add E$, and no non-trivial t-structure. This sharpens \cref{prop:local-exc} for this particular local algebra.
\item Let $(E_1,\dots,E_m)$ be a semi-orthogonal sequence of $\Pn{n_s}0$-like objects which is part of a Hom-dual triple. By \cref{thm:B}, it gives a standardly stratified algebra with $B_s\cong\kk[t]/(t^{n_s+1})$ if and only if neither sequence has negative extensions.
\item Conversely, by \cref{thm:C}, a standard module $\Delta(i)$ of a standardly stratified algebra with $B_i\cong\kk[t]/(t^{n+1})$ is $\Pn n0$-like.
\end{itemize}

The Hom-dual partners are of a different nature. Within its stratum, $N_s$ corresponds to the simple module over $B_s=\kk[t]/(t^{n_s+1})$. Its graded endomorphism algebra is $\Ext^\bullet_{B_s}(\kk,\kk)$, which by \cref{rem:Bonfert-Koszul} (with $k=0$) is as follows:
\begin{itemize}
\item for $n_s=1$, it is $\kk[s]$ with $\deg s=1$, so $N_s$ is ``$\Pn\infty1$-like'';
\item for $n_s\ge2$, its cohomology is $\kk[s,t]/(s^2)$ with $\deg s=1$ and $\deg t=2$, so $N_s$ is not $\mathbb{P}$-like.
\end{itemize}
The simplest example is $m=1$: then $E=B$ and $\Deltabar=N=\kk$ in $\Db(\modd B)$, and \cref{thm:B} recovers $\modd B$.


\subsection{The case \texorpdfstring{$k\ge1$}{k>=1}: local pairs from \texorpdfstring{$\mathbb{P}$}{P}-like objects}\label{sec:P-positive}

For $k\ge1$, a $\Pn nk$-like object has no negative self-extensions and satisfies $\End(P)=\kk$. It is therefore simple-minded in $\thick{P}$. We determine the heart of the associated t-structure, and the local pairs $(\Pi,P)$ in $\thick{P}$.


\begin{proposition}[{cf.\ \cite[\S3.2]{Bonf26}}]\label{prop:tP}
Let $P$ be $\Pn nk$-like with $k\ge1$. Then $\{P\}$ is a simple-minded collection in $\thick P$. Consequently $\thick P$ carries a bounded t-structure $t_P$ with heart $\cH_P:=\Filt(P)$, a length category with unique simple object $P$. The aisles of $t_P$ are described by \cref{lem:aisles} with $Q=S=P$.
\end{proposition}

\begin{proof}
Condition (SM1) holds because $\End^\bullet(P)$ is concentrated in degrees $0,k,\dots,nk\ge0$. Condition (SM2) holds because $\End(P)=\kk$ (as $k\ne0$), and (SM3) is clear. Now apply \cref{thm:smc}, and \cref{lem:aisles} with $Q=S=P$.
\end{proof}

\begin{remark}\label{rem:tP-unique}
The t-structure $t_P$ is the standard t-structure of the proper positive dg algebra $A_{n,k}$ \cite[Thm.~7.1]{KN13}, \cite[\S3.2]{Bonf26}. By \cref{thm:Bonfert}(2), every non-trivial t-structure on $\thick P$ is a shift of $t_P$. Hence, for strata generated by $\Pn nk$-like objects with $k\ge1$, the gluing data used in \cref{sec:gluing,sec:P-sequences} are determined up to shifts.
\end{remark}

\begin{proposition}[{cf.\ \cite[Lemmas~3.11 and~4.15, Cor.~4.16]{Bonf26}}]\label{prop:heart-P}
Let $P$ be $\Pn nk$-like with $k\ge1$ and $n\ge1$, and let $\cH_P$ be as in \cref{prop:tP}.
\begin{enumerate}
\item If $k\ge2$, then $\cH_P=\add P\simeq\modd\kk$.
\item If $k=n=1$, then $\cH_P$ is hereditary and has no non-zero projective objects. If $\thick P$ is algebraic, then $\real\colon\Db(\cH_P)\to\thick P$ is an equivalence.
\item Let $k=1$ and $n\ge2$. Let $t$ span $\Hom(P,P[1])$, and define $Q$ by the triangle
\[P\xrightarrow{a}Q\xrightarrow{b}P\xrightarrow{t}P[1].\]
Then:
\begin{itemize}
\item $Q\in\cH_P$ is a projective cover of $P$ in $\cH_P$, $\End(Q)\cong\kk[x]/(x^2)$, and $\cH_P\simeq\modd\kk[x]/(x^2)$;
\item $\Hom(Q,P)\cong\kk$, and for $l\ne0$ we have $\Hom(Q,P[l])\cong\kk$ if $l=n$ and $\Hom(Q,P[l])=0$ otherwise;
\item $\Hom(Q,Q[1])=0$.
\end{itemize}
\item In cases \textup{(1)} and \textup{(3)}, the comparison maps $\Ext^\bullet_{\cH_P}(P,P)\to\End^\bullet(P)$ are not bijective. In particular, no realisation functor $\Db(\cH_P)\to\thick P$ is an equivalence.
\end{enumerate}
\end{proposition}

\begin{proof}
Recall that $\Ext^1_{\cH_P}(X,Y)=\Hom(X,Y[1])$ for $X,Y\in\cH_P$ (\cref{prop:ext-comparison}).

(1) We have $\Ext^1_{\cH_P}(P,P)=\End^1(P)=0$ since $k\ge2$. Hence every object of $\Filt(P)$ is isomorphic to some $P^{\oplus r}$, and $\Hom(P,-)$ gives $\cH_P\simeq\modd\kk$.

(2) We have $\Ext^1_{\cH_P}(P,P)\cong\kk$. By \cref{prop:ext-comparison}, $\Ext^2_{\cH_P}(P,P)$ embeds into $\End^2(P)$, which is $0$ for $(k,n)=(1,1)$.
\begin{itemize}
\item By d\'evissage (long exact sequences in both variables), $\Ext^2_{\cH_P}$ vanishes identically. Since Yoneda extensions of length $\ge2$ factor through Yoneda products with $\Ext^2$, all $\Ext^{\ge2}_{\cH_P}$ vanish, so $\cH_P$ is hereditary.
\item Consequently $\chi(X,Y):=\dim\Hom(X,Y)-\dim\Ext^1(X,Y)$ is additive in both variables. It is therefore given by $\chi(X,Y)=\ell(X)\ell(Y)\chi(P,P)=0$, where $\ell$ denotes length.
\item If $R$ is projective, then $\dim\End(R)=\chi(R,R)=0$, so $R=0$.
\item The comparison maps $\Ext^l_{\cH_P}(P,P)\to\End^l(P)$ are bijective for $l\le1$ by \cref{prop:ext-comparison}, and both sides vanish for $l\ge2$. By d\'evissage and the five lemma, $\Ext^l_{\cH_P}(X,Y)\to\Hom(X,Y[l])$ is bijective for all $X,Y\in\cH_P$ and all $l$. Thus $\real$ is fully faithful, and it is an equivalence by \cref{thm:realisation}.
\end{itemize}

(3) Since $\cH_P$ is closed under extensions, $Q\in\cH_P$, and $0\to P\to Q\to P\to0$ is a non-split exact sequence in $\cH_P$. Let $W\in\thick P$. Applying $\Hom(-,W)$ to the triangle gives an exact sequence
\[\Hom(P[1],W)\xrightarrow{t^*}\Hom(P,W)\xrightarrow{b^*}\Hom(Q,W)\xrightarrow{a^*}\Hom(P,W)\xrightarrow{t[-1]^*}\Hom(P[-1],W).\]
For $W=P[l]$, the outer maps are, up to sign, multiplication by $t$, namely $\End^{l-1}(P)\to\End^l(P)$ and $\End^l(P)\to\End^{l+1}(P)$. Hence there is a short exact sequence
\[0\to\operatorname{coker}\big(\End^{l-1}(P)\xrightarrow{t}\End^l(P)\big)\to\Hom(Q,P[l])\to\ker\big(\End^l(P)\xrightarrow{t}\End^{l+1}(P)\big)\to0.\]
Since $\deg t=1$ and $t^n\ne0=t^{n+1}$:
\begin{itemize}
\item the cokernel is $\kk$ for $l=0$ and $0$ otherwise;
\item the kernel is $\kk$ (spanned by $t^n$) for $l=n$ and $0$ otherwise.
\end{itemize}
This proves the second item of (3); note that $0\neq n$.

Since $n\ge2$, we have $\Hom(Q,P[1])=0$, i.e.\ $\Ext^1_{\cH_P}(Q,P)=0$. As $P$ is the only simple object, d\'evissage gives $\Ext^1_{\cH_P}(Q,-)=0$, so $Q$ is projective. Since $\Hom(Q,P)\cong\kk$, the top of $Q$ is $P$, and $Q$ is a projective cover of $P$.

Applying $\Hom(Q,-)$ to the triangle gives an exact sequence
\[\Hom(Q,P[-1])\to\Hom(Q,P)\to\End(Q)\to\Hom(Q,P)\to\Hom(Q,P[1]).\]
Its outer terms vanish, so $\dim\End(Q)=2$. The algebra $\End(Q)$ is local because $Q$ has simple top; hence $\End(Q)\cong\kk[x]/(x^2)$. A Hom-finite length category with a projective generator $Q$ is equivalent to $\modd\End(Q)$ via $\Hom(Q,-)$, so $\cH_P\simeq\modd\kk[x]/(x^2)$. Finally, $\Hom(Q,Q[1])=\Ext^1_{\cH_P}(Q,Q)=0$.

(4) In case (1), $\Ext^{nk}_{\cH_P}(P,P)=0$ while $\End^{nk}(P)\ne0$. In case (3), $\Ext^{n+1}_{\kk[x]/(x^2)}(\kk,\kk)\cong\kk$, because the minimal projective resolution of $\kk$ is periodic, while $\End^{n+1}(P)=0$. A realisation functor induces the comparison maps on extension groups, so it cannot be fully faithful.
\end{proof}

\begin{remark}
Parts (1)--(3) overlap with \cite[Lemma~3.11]{Bonf26}, where the indecomposable objects of $\cH_P$ are listed and the spaces $\Hom(X_2,P[i])$ are computed; our $Q$ is Bonfert's $X_2$. They also overlap with \cite[Lemma~4.15]{Bonf26}, which shows that $\cH_P$ has an injective cogenerator unless $n=k=1$. The realisation statement in (2) corresponds to Bonfert's Koszul-dual model \cite[Cor.~4.16(2)]{Bonf26}, which identifies $\thick P$ with the thick subcategory of $\mathcal{D}(\kk[[t]])$ generated by the simple module. What is new in \cref{prop:heart-P} is only the explicit identification $\cH_P\simeq\modd\kk[x]/(x^2)$ in (3) and the reformulation of these facts in terms of local pairs and their Hom-duality (\cref{cor:P-local-pairs}), which is what the gluing results use. We have given the proofs in the form needed below.
\end{remark}

\begin{remark}\label{rem:koszul}
\leavevmode
\begin{enumerate}
\item For $n=1$, parts (1) and (2) are consistent with \cref{thm:HJY}. For $w=1$, $\thick P$ is the derived category of finite-dimensional nilpotent $\kk[x]$-modules; cf.\ \cite[Lemma~3.11(3), Cor.~4.16(2)]{Bonf26}. For $w=2$, part~(1) is exactly the phenomenon of \cref{ex:BB-5.11}.
\item The simple module $\kk$ generates $\Db(\modd\kk[x]/(x^2))$, and $\End^\bullet(\kk)\cong\kk[s]$ with $\deg s=1$. So $\kk$ is a ``$\Pn\infty1$-like'' object (\cref{rem:Bonfert-Koszul} with $n=1$, $k=0$). A $\Pn n1$-like object with $n\ge2$ has the same heart, but its endomorphism algebra is truncated at $t^{n+1}=0$. Meanwhile, the regular module $\kk[x]/(x^2)$ is $\Pn10$-like. Thus the cases $k=0$ and $k=1$ are Koszul dual to each other, which is the mechanism behind \cref{thm:HJY}.
\item The non-vanishing $\Hom(Q,P[n])\ne0$ in (3) and $\End^{k}(P)\ne0$ in (1) are exactly the obstructions to full Hom-duality of the corresponding local pairs (\cref{def:local-pair}). The vanishing of $\Hom(Q,P[2])$ (i.e.\ $n\ne2$) and of $\End^2(P)$ (i.e.\ $k\ne2$) is Hom-duality up to degree $2$, which is the hypothesis of \cref{thm:A}.
\end{enumerate}
\end{remark}



\begin{corollary}\label{cor:P-local-pairs}
Let $P$ be $\Pn nk$-like with $k\ge1$ and $n\ge1$.
\begin{enumerate}
\item If $k\ge2$, then $(P,P)$ is a local pair in $\thick{P}$ with Hom-duality up to degree $k-1$, but not up to degree $k$.
\item If $k=1$ and $n\ge2$, then $(Q,P)$ is a local pair in $\thick{P}$ with Hom-duality up to degree $n-1$, but not up to degree $n$.
\item If $k=n=1$, there is no local pair of the form $(\Pi,P)$ in $\thick{P}$.
\end{enumerate}
\end{corollary}

\begin{proof}
(LP1) is \cref{prop:tP}. For (1), $\End^\bullet(P)$ is concentrated in the degrees $0,k,\dots,nk$. For (2), use \cref{prop:heart-P}(3), and note that $Q\in\cH_P$ has no negative self-extensions. For (3), if $(\Pi,P)$ were a local pair, then $\Pi$ would be a non-zero projective object of $\cH_P$ by \cref{prop:local-pair}(2), contradicting \cref{prop:heart-P}(2).
\end{proof}


\subsection{Semi-orthogonal sequences of \texorpdfstring{$\mathbb{P}$}{P}-like objects}\label{sec:P-sequences}

\begin{definition}\label{def:P-sequence}
A \emph{$\mathbb{P}$-sequence} in $\cD$ is a sequence $(P_1,\dots,P_m)$ with the following properties:
\begin{itemize}
\item each $P_s$ is $\Pn{n_s}{k_s}$-like for some $n_s\ge0$ and $k_s\ge1$ (for $n_s=0$ the value of $k_s$ is irrelevant);
\item $\Hom^\bullet(P_t,P_s)=0$ for $t>s$.
\end{itemize}
It is \emph{full} if $\thick{P_1,\dots,P_m}=\cD$. Put $\cC_s:=\thick{P_1,\dots,P_s}$. A \emph{left dual} is a sequence $(F_m,\dots,F_1)$ together with morphisms $\varphi_s\colon P_s\to F_s$ such that, for every $s$, the cocone of $\varphi_s$ lies in $\cC_{s-1}$, and $\Hom^\bullet(\cC_{s-1},F_s)=0$.
\end{definition}

The left dual is the mirror image of the right dual in \cref{def:triple}. By the argument of \cref{lem:right-dual}, it is unique up to unique isomorphism. It exists if and only if each $\cC_{s-1}$ is right admissible in $\cC_s$, and then $F_s$ is the left mutation of $P_s$ through $\cC_{s-1}$; compare \cref{rem:HD4}. For exceptional sequences this is the usual left dual, and it always exists.

\begin{lemma}\label{lem:P-triple}
Let $(P_1,\dots,P_m)$ be a full $\mathbb{P}$-sequence with a left dual $(F_m,\dots,F_1)$.
\begin{enumerate}
\item $\cC_s=\thick{F_1,\dots,F_s}$ for all $s$, and $\End^\bullet(F_s)\cong\End^\bullet(P_s)$. The filtration $(\cC_s)$ is admissible.
\item Define $\Delta_s$ as follows. If $n_s=0$ or $k_s\ge2$, put $\Delta_s:=P_s$. If $k_s=1$ and $n_s\ge2$, let $\Delta_s:=Q_s$ be the cocone of a non-zero morphism $P_s\to P_s[1]$, computed in $\cD$. Assume that no $(k_s,n_s)$ equals $(1,1)$. Then $(\Delta_\bullet,P_\bullet,F_\bullet)$, with $\delta_s=\varphi_s$, is a stratifying triple in $\cD$. Its local pairs correspond to $(\Delta_s,P_s)$ under $\cT_s\simeq\thick{P_s}$. It has Hom-duality up to degree $\min_s d_s$, where
\[d_s=\infty\ \text{ if } n_s=0,\qquad d_s=k_s-1\ \text{ if } k_s\ge2,\qquad d_s=n_s-1\ \text{ if } k_s=1.\]
\end{enumerate}
\end{lemma}

\begin{proof}
(1) Since $\cone(\varphi_s)\in\cC_{s-1}$, we have $F_s\in\cC_s$ and $\cC_s=\thick{\cC_{s-1}\cup\{F_s\}}$. Induction gives $\cC_s=\thick{F_1,\dots,F_s}$.

Semi-orthogonality of the $P$'s gives $\Hom^\bullet(P_s,\cC_{s-1})=0$. So, applying $\Hom(P_s,-)$ and $\Hom(-,F_s)$ to the triangle relating $P_s$ and $F_s$ gives
\[\End^\bullet(F_s)\cong\Hom^\bullet(P_s,F_s)\cong\End^\bullet(P_s).\]
Finally, \cref{lem:sod}(2) gives both adjoints of $\cC_{s-1}\subset\cC_s$:
\begin{itemize}
\item the pair $(\thick{P_s},\cC_{s-1})$ with $\Hom^\bullet(P_s,\cC_{s-1})=0$ gives the left adjoint;
\item the pair $(\cC_{s-1},\thick{F_s})$ with $\Hom^\bullet(\cC_{s-1},F_s)=0$ gives the right adjoint.
\end{itemize}

(2) In both cases $\Delta_s\in\thick{P_s}\subset{}^\perp\cC_{s-1}\cap\cC_s$. We check the conditions of \cref{def:triple}.
\begin{itemize}
\item \emph{(ST2) and (ST3)} follow from (1), the semi-orthogonality of the $F$'s (as $F_s\in\cC_s\subset\cC_{t-1}$ for $s<t$) and $\End^\bullet(P_s)\cong\kk[t]/(t^{n_s+1})$ with $k_s\ge1$.
\item \emph{(ST4)} holds with $\Deltabar_s=P_s$: indeed $\cone(\varphi_s)\in\cC_{s-1}$, and $\Hom^\bullet(P_s,F_t)=0$ for $t<s$ because $F_t\in\cC_{s-1}$.
\item \emph{(ST1).} For $t<s$, $\Hom^\bullet(\Delta_s,F_t)=0$ because $F_t\in\cC_{s-1}$. For $t>s$, it vanishes because $\Delta_s\in\cC_{t-1}$. For $t=s$, applying $\Hom(\Delta_s,-)$ to the triangle relating $P_s$ and $F_s$ gives $\Hom^\bullet(\Delta_s,F_s)\cong\Hom^\bullet(\Delta_s,P_s)$, which is computed in $\thick{P_s}$. By \cref{cor:P-local-pairs}, $(\Delta_s,P_s)$ is a local pair in $\thick{P_s}$ with Hom-duality up to degree $d_s$; for $n_s=0$ it is the pair $(\kk,\kk)$.
\item \emph{(ST5)} holds for the same reason.
\end{itemize}
Finally, by \cref{lem:peel}(1) the stratum $\cT_s$ is identified with ${}^\perp\cC_{s-1}\cap\cC_s=\thick{P_s}$, and under this identification $j_{s!}$ is the inclusion.
\end{proof}

\begin{theorem}\label{thm:Pseq}
Let $(P_1,\dots,P_m)$ be a full $\mathbb{P}$-sequence with a left dual $(F_m,\dots,F_1)$, and suppose that for every $s$ one of the following holds:
\[n_s=0;\qquad k_s\ge3;\qquad k_s=1\text{ and }n_s\ge3.\]
Suppose moreover that $\Hom(P_s,P_t[l])=0$ and $\Hom(F_s,F_t[l])=0$ for all $s,t$ and all $l<0$. Let $\Delta_s$ be as in \cref{lem:P-triple}, let $G$ be the iterated universal extension of $(\Delta_1,\dots,\Delta_m)$, and let $\cA$ be the heart of the t-structure of the triple $(\Delta_\bullet,P_\bullet,F_\bullet)$.
\begin{enumerate}
\item $\cA\simeq\modd\End(G)$ is standardly stratified with respect to $1<\dots<m$. Its standard objects are $\Delta_s$, its proper standard objects are $P_s$, and its proper costandard objects are $F_s$. We have $B_s\cong\kk$ if $n_s=0$ or $k_s\ge3$, and $B_s\cong\kk[x]/(x^2)$ if $k_s=1$. If $k_s=1$, there is an exact sequence $0\to P_s\to\Delta_s\to P_s\to0$ in $\cA$.
\item If $n_s=0$ or $k_s\ge3$ for all $s$, then $\cA$ is a highest weight category with standard objects $P_s$ and costandard objects $F_s$.
\item If $\cD$ is algebraic, then $\real\colon\Db(\cA)\to\cD$ is an equivalence if and only if $n_s=0$ for all $s$.
\end{enumerate}
The t-structure of the triple is the t-structure glued from the t-structures $t_{P_s}$ on the strata $\thick{P_s}$. By \cref{rem:tP-unique}, every t-structure glued from non-trivial t-structures on the strata is obtained from it by shifting the strata, i.e.\ by replacing $P_s$ by $P_s[p(s)]$ for some $p(s)\in\mathbb{Z}$.
\end{theorem}

\begin{proof}
By \cref{lem:P-triple}, $(\Delta_\bullet,P_\bullet,F_\bullet)$ is a stratifying triple with Hom-duality up to degree $\min_sd_s\ge2$. Since $\Delta_s\in\Filt(P_s)$ (for $k_s=1$ this uses the defining triangle of $Q_s$), the vanishing of negative extensions among the $P_s$ implies it among the $\Delta_s$. So (NH) holds, and \cref{thm:A} applies. This gives (1), except for the statements about $P_s$, and also (3): full Hom-duality holds if and only if all $d_s=\infty$, i.e.\ all $n_s=0$.

\emph{The proper standard objects.} By \cref{prop:triple-t}(3), $P_s\in\cA$. If $k_s=1$, the triangle $P_s\to\Delta_s\to P_s\to P_s[1]$ has all its vertices in $\cA$, so it gives the exact sequence in~(1). We have $[\Delta_s:L_s]=2$, since $j_s^*\Delta_s=Q_s$ has length two in $\cH_{P_s}$, and $[P_s:L_s]=1$.

Let $N_0$ be the kernel of $\Delta_s\twoheadrightarrow\Deltabar(s)$. Recall that $\Deltabar(s)$ is the largest quotient of $\Delta_s$ with a single composition factor $L_s$; its other composition factors lie in $\cA_{s-1}$.
\begin{itemize}
\item The quotient $\Delta_s/P_s\cong P_s$ has this property, so $N_0\subset P_s$.
\item Since $[\Delta_s:L_s]=2$ and $[\Deltabar(s):L_s]=1$, we have $[N_0:L_s]=1=[P_s:L_s]$. So $P_s/N_0$ has no composition factor $L_s$, i.e.\ $P_s/N_0\in\cA_{s-1}$.
\item But $P_s={}^p\!j_{s!}S_s$ has no non-zero quotient in $\cA_{s-1}$ (\cref{prop:glued-heart}(2)).
\end{itemize}
Hence $N_0=P_s$, and $\Deltabar(s)\cong\Delta_s/P_s\cong P_s$. If $n_s=0$ or $k_s\ge3$, then $\Delta_s=P_s$ and $[P_s:L_s]=1$, so $\Deltabar(s)=\Delta(s)=P_s$.

(2) follows from (1), since standardly stratified with all $B_s=\kk$ means quasi-hereditary, and then $\nabla=\nablabar$.

The last statement follows from \cref{prop:triple-t}(1) and \cref{rem:tP-unique}.
\end{proof}

\begin{proposition}\label{prop:P11}
Let $(P_1,\dots,P_m)$ be a full $\mathbb{P}$-sequence with a left dual, and suppose $(k_s,n_s)=(1,1)$ for some $s$. Glue the t-structures $t_{P_t}$ along the admissible filtration $(\cC_t)$ of \cref{lem:P-triple}(1). If $F_s$ lies in the glued heart $\cA$, then the simple object $L_s=j_{s!*}P_s$ has no projective cover in $\cA$. In particular $\cA$ is not equivalent to $\modd\Lambda$ for any finite-dimensional algebra $\Lambda$.
\end{proposition}

\begin{proof}
Suppose $R$ is a projective cover of $L_s$ in $\cA$, and let $\bar R$ be its largest quotient lying in $\cA_s$; then $\bar R$ is a projective cover of $L_s$ in $\cA_s$.
\begin{itemize}
\item For $M\in\cH_{P_s}$ we have $j_{s*}M\in\Filt(F_s)\subset\cA_s$. Hence
\[\Ext^1_{\cH_{P_s}}(j_s^*\bar R,M)=\Hom(\bar R,j_{s*}M[1])=\Ext^1_{\cA_s}(\bar R,j_{s*}M)=0,\]
so $j_s^*\bar R$ is projective.
\item $j_s^*\bar R\ne0$: otherwise $\bar R\in\cA_{s-1}$, and then $\Hom(\bar R,L_s)=0$ because $L_s$ has no non-zero subobject in $\cA_{s-1}$.
\end{itemize}
This contradicts \cref{prop:heart-P}(2).
\end{proof}

\begin{example}[A case with $k=2$]\label{ex:k2}
Let $\Gamma$ be the graded $\kk$-linear category (a dg category with zero differential) with two objects $1,2$ and the following morphism spaces:
\begin{itemize}
\item $\Gamma(1,1)=\kk$ and $\Gamma(2,2)=\kk[t]/(t^2)$ with $\deg t=2$;
\item $\Gamma(1,2)=\kk f$ with $\deg f=0$, and $\Gamma(2,1)=0$;
\item $t\circ f=0$.
\end{itemize}
Let $\cD=\operatorname{perf}(\Gamma)$, and let $P_1,P_2$ be the representable objects, so that $\Hom^\bullet(P_i,P_j)=\Gamma(i,j)$. Then $(P_1,P_2)$ is a full $\mathbb{P}$-sequence with $n_1=0$ and $P_2$ a $\Pn12$-like object.

It has a left dual: $F_1=P_1$, and $F_2=\cone(P_1\xrightarrow{f}P_2)$ with $\varphi_2$ the canonical morphism. Here $\Hom^\bullet(P_1,F_2)=0$ because $\Hom^\bullet(P_1,P_2)=\kk f$. The triple $(P_\bullet,P_\bullet,F_\bullet)$ has Hom-duality only up to degree $1$, so \cref{thm:Pseq} does not apply.

Nevertheless the conclusion of \cref{thm:Pseq}(2) holds. Negative extensions vanish:
\begin{itemize}
\item $\Hom^\bullet(P_1,P_2)=\kk f$ is concentrated in degree $0$;
\item $\Hom^\bullet(F_1,F_2)=0$;
\item $\Hom^l(F_2,F_1)\cong\End^{l-1}(P_1)$ is concentrated in degree $1$;
\item $\End^\bullet(F_2)\cong\End^\bullet(P_2)$.
\end{itemize}
So $P_s,F_s\in\cA$ by \cref{prop:triple-t}(3). The triangle $P_1\to P_2\to F_2\to P_1[1]$ gives an exact sequence $0\to P_1\to P_2\to F_2\to0$ in $\cA$. Hence the simple objects are $L_1=P_1$ and $L_2=F_2$. Their extensions are:
\begin{itemize}
\item $\Ext^1_\cA(L_2,L_1)=\Hom(F_2,P_1[1])\cong\kk$;
\item $\Ext^1_\cA(L_1,L_2)=\Hom(P_1,F_2[1])=0$;
\item $\Ext^1_\cA(L_1,L_1)=\Ext^1_\cA(L_2,L_2)=0$.
\end{itemize}
Moreover $\Ext^1_\cA(P_1,-)$ and $\Ext^1_\cA(P_2,-)$ vanish on both simple objects, so $P_1,P_2$ are the indecomposable projective objects. Thus $\cA\simeq\modd\kk\vec A_2$ is highest weight, with standard objects $P_1,P_2$ and costandard objects $F_1,F_2$.

This is consistent with \cref{rem:star}: here $V_1=\Ext^1_\cA(P_1,P_2)=0$, so Hom-duality in degree $2$ is not needed for the top stratum. The realisation functor is not an equivalence, since $\Hom(P_2,P_2[2])\ne0$ while $\Ext^2_{\kk\vec A_2}=0$.
\end{example}

\begin{remark}[The obstruction for $k=2$]\label{rem:obstruction}
Suppose all strata have semisimple hearts, i.e.\ $k_s\ge2$ or $n_s=0$. In the induction step (\cref{prop:step}(2)), the only point where Hom-duality in degree $2$ is needed for an earlier stratum $t<m$ is the case $l=1$ (step (e)). There the kernel $K$ of $\Delta_m\twoheadrightarrow L_m$ lies in $\cA_{m-1}$, and one needs $\Hom(G_s,K[2])=0$.

If some $k_t=2$ with $t<m$, this group can be non-zero. Whether $\Ext^1_\cA(G_s,L_m)$ vanishes then depends on the connecting map $\Hom(G_s,L_m[1])\to\Hom(G_s,K[2])$. This map is determined by higher (Massey-type) composition data, and not by the graded endomorphism algebras alone; compare the non-formality of the Koszul duals in \cref{rem:Bonfert-Koszul}.
\end{remark}

\begin{question}\label{q:k2}
Let $(P_1,\dots,P_m)$ be a full $\mathbb{P}$-sequence with a left dual, without negative extensions among the $P$'s and among the $F$'s, and suppose that some $k_s=2$ or some $(k_s,n_s)=(1,2)$. Is the heart of the t-structure glued from the $t_{P_s}$ Deligne finite? If so, is it standardly stratified with standard objects $\Delta_s$, as in \cref{thm:Pseq}?
\end{question}

Natural test cases come from geometry:
\begin{itemize}
\item the structure sheaves of hyperk\"ahler manifolds are $\Pn n2$-objects;
\item by \cite[Thm.~B]{BC25}, as quoted in \cite{Bonf26}, all indecomposable perverse sheaves on $\mathbb{P}^n$ are $\mathbb{P}[2]$-like, and most of them are not Calabi--Yau.
\end{itemize}
It would be interesting to know whether such objects can be arranged into $\mathbb{P}$-sequences with left duals satisfying the hypotheses above.

% ===============================================================
\section{Further questions}\label{sec:questions}
% ===============================================================

\begin{question}[Silting gluing]\label{q:silting}
Can \cref{thm:B} be derived from the theory of gluing silting objects \cite{LVY14}?

On the perfect part $\thick{E_1,\dots,E_m}$, only the adjoints $i^*$ and $j_!$ are available in general, and they are given by the $E_s$. On $\cD$, the adjoints $i^!$ and $j_*$ are automatic (\cref{rem:HD4}), and they are given by the $N_s$. One would like to glue silting objects on the perfect side along these ``half recollements'', glue the simple-minded objects $N_s$ on the other side, and use the Hom-duality (HD1) as the compatibility condition, in the spirit of Koenig--Yang (\cref{thm:KY}).
\end{question}

\begin{question}[Without a right dual]
Is there a version of \cref{thm:A,thm:B} assuming only (ST1)--(ST3) and (ST5), i.e.\ only the existence of the functors $j_{s*}$, which is automatic by \cref{lem:sod}? This would require a one-sided gluing of t-structures along semi-orthogonal decompositions (cf.\ \cite{CP10}). It would remove the only hypothesis of \cref{thm:B} that is not automatic, and might answer \cref{q:HD4}.
\end{question}

\begin{question}[Classification of glued t-structures]\label{q:classification}
For a full $\mathbb{P}$-sequence with left dual and all $k_s\ge1$, \cref{thm:Bonfert} implies that the t-structures on $\cD$ glued from non-trivial t-structures on the strata are parametrised by perversity functions $p\colon\{1,\dots,m\}\to\mathbb{Z}$. Which t-structures on $\cD$ are glued in this sense? For which perversities do the objects $P_s[p(s)]$ and $F_s[p(s)]$ lie in the heart, as for the middle perversity in \cite[\S5.1]{BB26}?
\end{question}

\begin{question}[Ringel duality]
By \cite[Thm.~3.2]{BB26}, the Ringel dual of a highest weight category glued along an exceptional sequence is glued along the left dual sequence. Characteristic tilting modules and Ringel duality are available for standardly and properly stratified algebras (see e.g.\ \cite{AHLU00}). Is there a ``dual'' stratifying triple, built from $(E_\bullet,\Deltabar_\bullet,N_\bullet)$ in \cref{thm:B}, whose associated heart is the Ringel dual?
\end{question}

% ===============================================================
\begin{thebibliography}{BBD82}
% ===============================================================

\bibitem[AHLU00]{AHLU00}
I.~\'Agoston, D.~Happel, E.~Luk\'acs, L.~Unger, \emph{Standardly stratified algebras and tilting}, J.~Algebra \textbf{226} (2000), 144--160.

\bibitem[AI12]{AI12}
T.~Aihara, O.~Iyama, \emph{Silting mutation in triangulated categories}, J.~Lond. Math. Soc. (2) \textbf{85} (2012), 633--668.

\bibitem[AN09]{AN09}
S.~Al-Nofayee, \emph{Simple objects in the heart of a t-structure}, J.~Pure Appl. Algebra \textbf{213} (2009), 54--59.

\bibitem[BB26]{BB26}
A.~Bodzenta, A.~Bondal, \emph{Highest weight categories via pairs of dual exceptional sequences}, arXiv:2601.22004 (2026).

\bibitem[BBD82]{BBD82}
A.~A.~Beilinson, J.~Bernstein, P.~Deligne, \emph{Faisceaux pervers}, Ast\'erisque \textbf{100} (1982).

\bibitem[BC25]{BC25}
L.~Bonfert, A.~Cipriani, \emph{Serre functor and $\mathbb{P}$-objects for perverse sheaves on $\mathbb{P}^n$}, arXiv:2506.06051 (2025).

\bibitem[Bei87]{Bei87}
A.~A.~Beilinson, \emph{On the derived category of perverse sheaves}, in: K-theory, arithmetic and geometry (Moscow, 1984--1986), Lecture Notes in Math. \textbf{1289}, Springer, 1987, 27--41.

\bibitem[BK89]{BK89}
A.~I.~Bondal, M.~M.~Kapranov, \emph{Representable functors, Serre functors, and reconstructions}, Izv. Akad. Nauk SSSR Ser. Mat. \textbf{53} (1989), 1183--1205.

\bibitem[Bon89]{Bon89}
A.~I.~Bondal, \emph{Representations of associative algebras and coherent sheaves}, Izv. Akad. Nauk SSSR Ser. Mat. \textbf{53} (1989), 25--44.

\bibitem[Bon10]{Bon10}
M.~V.~Bondarko, \emph{Weight structures vs.\ t-structures; weight filtrations, spectral sequences, and complexes (for motives and in general)}, J.~K-Theory \textbf{6} (2010), 387--504.

\bibitem[Bonf26]{Bonf26}
L.~Bonfert, \emph{Weight structures and t-structures on the triangulated category generated by a $\mathbb{P}$-object}, arXiv:2610.21859 (2026).

\bibitem[CP10]{CP10}
J.~Collins, A.~Polishchuk, \emph{Gluing stability conditions}, Adv. Theor. Math. Phys. \textbf{14} (2010), 563--607.

\bibitem[CPS88]{CPS88}
E.~Cline, B.~Parshall, L.~Scott, \emph{Finite-dimensional algebras and highest weight categories}, J.~Reine Angew. Math. \textbf{391} (1988), 85--99.

\bibitem[CPS96]{CPS96}
E.~Cline, B.~Parshall, L.~Scott, \emph{Stratifying endomorphism algebras}, Mem. Amer. Math. Soc. \textbf{124} (1996), no.~591.

\bibitem[Dla96]{Dla96}
V.~Dlab, \emph{Quasi-hereditary algebras revisited}, An. \c{S}tiin\c{t}. Univ. Ovidius Constan\c{t}a Ser. Mat. \textbf{4} (1996), 43--54.

\bibitem[DR92]{DR92}
V.~Dlab, C.~M.~Ringel, \emph{The module theoretical approach to quasi-hereditary algebras}, in: Representations of algebras and related topics (Kyoto, 1990), London Math. Soc. Lecture Note Ser. \textbf{168}, Cambridge Univ. Press, 1992, 200--224.

\bibitem[ES03]{ES03}
K.~Erdmann, C.~S\'aenz, \emph{On standardly stratified algebras}, Comm. Algebra \textbf{31} (2003), 3429--3446.

\bibitem[HJY13]{HJY13}
T.~Holm, P.~J\o rgensen, D.~Yang, \emph{Sparseness of t-structures and negative Calabi--Yau dimension in triangulated categories generated by a spherical object}, Bull. Lond. Math. Soc. \textbf{45} (2013), 120--130.

\bibitem[HK19]{HK19}
A.~Hochenegger, A.~Krug, \emph{Formality of $\mathbb{P}$-objects}, Compos. Math. \textbf{155} (2019), no.~5.

\bibitem[HT06]{HT06}
D.~Huybrechts, R.~Thomas, \emph{$\mathbb{P}$-objects and autoequivalences of derived categories}, Math. Res. Lett. \textbf{13} (2006), 87--98.

\bibitem[Kal17]{Kal17}
M.~Kalck, \emph{Derived categories of quasi-hereditary algebras and their derived composition series}, in: Representation theory---current trends and perspectives, EMS Ser. Congr. Rep., Eur. Math. Soc., 2017, 269--308.

\bibitem[Kel94]{Kel94}
B.~Keller, \emph{Deriving DG categories}, Ann. Sci. \'Ec. Norm. Sup\'er. (4) \textbf{27} (1994), 63--102.

\bibitem[KN13]{KN13}
B.~Keller, P.~Nicol\'as, \emph{Weight structures and simple dg modules for positive dg algebras}, Int. Math. Res. Not. IMRN \textbf{2013}, no.~5.

\bibitem[Kru18]{Kru18}
A.~Krug, \emph{Varieties with $\mathbb{P}$-units}, Trans. Amer. Math. Soc. \textbf{370} (2018), no.~11.

\bibitem[KY14]{KY14}
S.~Koenig, D.~Yang, \emph{Silting objects, simple-minded collections, t-structures and co-t-structures for finite-dimensional algebras}, Doc. Math. \textbf{19} (2014), 403--438.

\bibitem[Let24]{Let24}
J.~C.~Letz, \emph{Realization functors in algebraic triangulated categories}, arXiv:2412.07455 (2024).

\bibitem[LVY14]{LVY14}
Q.~Liu, J.~Vit\'oria, D.~Yang, \emph{Gluing silting objects}, Nagoya Math. J. \textbf{216} (2014), 117--151.

\bibitem[Maz04]{Maz04}
V.~Mazorchuk, \emph{On finitistic dimension of stratified algebras}, Algebra Discrete Math. (2004), no.~3, 77--88.

\bibitem[MMS04]{MMS04}
E.~N.~Marcos, O.~Mendoza, C.~S\'aenz, \emph{Stratifying systems via relative simple modules}, J.~Algebra \textbf{280} (2004), 472--487.

\bibitem[Pau08]{Pau08}
D.~Pauksztello, \emph{Compact corigid objects in triangulated categories and co-t-structures}, Cent. Eur. J. Math. \textbf{6} (2008), 25--42.

\bibitem[ST01]{ST01}
P.~Seidel, R.~Thomas, \emph{Braid group actions on derived categories of coherent sheaves}, Duke Math. J. \textbf{108} (2001), 37--108.

\end{thebibliography}

\end{document}
